\chapter{実験の補遺}
\label{sec:supplement}

ここでは章ごとに、主な主張とその根拠をまとめる。根拠とは、それを測定した実験、何を測定したか、どこで発表されたかである。右の列の状態は、主張がどれだけ確かかを示す。\emph{測定済み}は直接の実験があること、\emph{理論}は本文や古典的な研究で導かれた数学的帰結、\emph{モデル}は本書のモデルによる計算結果（スクリプトは \texttt{models/} フォルダにある）、\emph{推定}は直接の測定のない桁の見積もり、\emph{仮説}はもっともらしいがまだ実験で証明されていない説明である。実験がない場合は、そう明記した。

\section{第\ref{sec:neurontypes}章　ニューロンの種類}

この章の数値は、ヒトの脳で細胞の直接計数があるところではそれに基づく。「ニューロン1個に神経伝達物質1種類」という規則は、細胞アトラスと、例外も見つけた実験に基づく。\nt{DOPA}の役割はスパイクの記録に基づき、その他のブロードキャストチャネルの役割は仮説に基づく。

{\footnotesize
\setlength{\tabcolsep}{3pt}\renewcommand{\arraystretch}{1.15}
\begin{longtable}{>{\raggedright}p{0.5cm}>{\raggedright}p{4.0cm}>{\raggedright}p{5.6cm}>{\raggedright}p{2.3cm}>{\raggedright\arraybackslash}p{1.7cm}}
\toprule
No. & 主張 & 実験と測定されたもの & 出典 & 状態 \\
\midrule\endfirsthead
\toprule
No. & 主張 & 実験と測定されたもの & 出典 & 状態 \\
\midrule\endhead
1 & ヒトの脳のニューロンは約$8.6\cdot10^{10}$個で、\nt{GLUT}ニューロンのほとんどは小脳の小さなニューロンである（表~\ref{tab:types}） & 等方性分画法：組織を溶かして核の均一な懸濁液にし、ニューロンマーカーNeuNをもつ核を数える。成人男性のニューロンは$(86.1\pm8.1)\cdot10^9$個。そのうち大脳皮質にあるのは$19\,\%$にすぎず、大部分は小脳にある & \cite{Azevedo2009} & 測定済み \\
2 & ほとんどすべてのニューロンは主要な神経伝達物質を1種類もつ（命題~\ref{prop:dale}）が、例外もある & マウス脳のトランスクリプトーム・アトラス：約$4\cdot10^6$個の細胞、$5322$クラスタ。合成と小胞充填の遺伝子から各クラスタに神経伝達物質を割り当てた。例外は直接測定されている。スライス標本の\nt{DOPA}ニューロンは同じ小胞トランスポーターを通じてGABAも放出し、線条体のニューロンを抑制する。\nt{DOPA}軸索の一部では、グルタミン酸とドーパミンが同じ配線の異なる終末から出る（電子顕微鏡と光遺伝学） & \cite{Strata1999,Yao2023,Tritsch2012,Zhang2015} & 測定済み \\
3 & 重み行列の列全体が同じ符号をもつ~\eqref{eq:dalesign} & 命題~\ref{prop:dale}と、グルタミン酸受容体はほぼすべて興奮性でGABA受容体は抑制性であることから章内で導出。「1個のニューロンのすべての受け手が同じ符号の重みを受け取る」ことを一般規則として直接検証した例はない。反例は\nt{DOPA}ニューロンによるGABAの共放出（2行目）と、異なる符号の受容体をもつ受け手（11行目） & 章内の導出 & 理論 \\
4 & 大脳皮質のニューロンの4分の3から5分の4が\nt{GLUT}、5分の1から4分の1が\nt{GABA} & サル大脳皮質の$10$領域で、皮質の全層を貫く幅$50$\,\textmu mの柱状領域をGABAで免疫染色。$8$領域ではGABAニューロンが全ニューロンの約$25\,\%$、一次視覚野では$20\,\%$未満 & \cite{Hendry1987} & 測定済み \\
5 & \nt{GLUT}ニューロンの出力は約$10^4$個 & 剖検脳の立体解析学：若年男性の新皮質にシナプス$1.64\cdot10^{14}$個（電子顕微鏡、ダイセクター法）、ニューロン約$2.3\cdot10^{10}$個。比をとるとニューロン1個あたり$\sim7\cdot10^3$シナプス。平均すれば入力数は出力数に等しい & \cite{Tang2001synapses,Pakkenberg1997} & 測定済み \\
6 & \nt{DOPA}ニューロンの出力は$10^5$--$10^6$個の終末に枝分かれする。これは標的の被覆範囲からの推定であって計数ではない & ラットの個々の\nt{DOPA}ニューロンの膜をウイルスで標識し、軸索を完全に再構成：1個のニューロンが線条体の体積の$0.45$から$5.7\,\%$（平均$2.7\,\%$）を覆う。測定したのは被覆範囲で、終末数ではない。終末数は被覆範囲からオーダーとして導いたもので、終末の直接計数はない。\nt{SERO}と\nt{NORA}の終末数は粗い推定 & \cite{Matsuda2009} & 測定済み \\
7 & ブロードキャスト型ニューロンは少ない：\nt{DOPA} $\sim5\cdot10^5$（二つの群のうち主要なほう、下限）、\nt{SERO} $\sim3\cdot10^5$（二つの部分計数の和）、\nt{HIST} $\sim6\cdot10^4$（表~\ref{tab:types}、図~\ref{fig:typepie}） & ヒトでの立体解析学的計数：黒質の色素含有ニューロン$550\,000$個（腹側被蓋野を除く）。背側縫線核のニューロン$235\,000$個（核の全ニューロン）と、橋被蓋のセロトニンニューロンさらに$125\,000$個。視床下部のヒスタミンニューロン約$64\,000$個。\nt{NORA}、\nt{GLYC}、\nt{ACET}、\nt{ADRE}には直接の計数がなく、表の値は粗いオーダー推定である & \cite{Pakkenberg1991,Baker1990,Baker1991,Airaksinen1991} & 測定済み \\
8 & 隙間のグルタミン酸は約1ミリモーラーまで跳ね上がり、$\sim1\ms$で減衰する & シナプス伝達中に、速く解離する競合的阻害薬がNMDA受容体から追い出される程度から求めた（海馬ニューロンの培養）：ピーク$1.1$\,mM、減衰時定数$1.2\ms$ & \cite{Clements1992} & 測定済み \\
9 & 動作中の皮質では興奮性電流と抑制性電流がほぼ釣り合い、ニューロンは閾値付近にいる & 理論：強い結合をもつネットワークは自ら釣り合いの状態に入る。実験：フェレット前頭前皮質のin vivo記録で、「アップ」状態の間、コンダクタンスが$\sim21$\,nS変化してもシナプス電流の反転電位は数百ミリ秒にわたり約$-37$\,mVに保たれる。興奮と抑制が一緒に増減する & \cite{vanVreeswijk1996,Haider2006} & 測定済み \\
10 & \nt{DOPA}ニューロンは報酬予測誤差~\eqref{eq:rpe}を伝える（図~\ref{fig:rpe}） & サルの\nt{DOPA}ニューロンのスパイク：予告なしのジュースにバースト。学習後は合図にバーストし、予測されたジュースには応答しない。約束されたジュースが与えられないと発火率が落ち込む。定量的な実験では、発火率は現在の報酬と過去の報酬の加重平均との差に比例するが、それは正の誤差に限られる。式~\eqref{eq:rpe}そのものは時間差分アルゴリズムである & \cite{Schultz1997,Bayer2005,SuttonBarto2018} & 測定済み \\
11 & 符号と速さを決めるのは受容体：イオンチャネル型は1ミリ秒未満、代謝型ははるかに遅い（命題~\ref{prop:receptor}） & 海馬ニューロンの膜パッチに短いグルタミン酸パルスを与えると、イオンチャネル型受容体の電流は$0.2$--$0.6\ms$で立ち上がり（$20$から$80\,\%$）、$2$--$3\ms$で減衰する。錐体ニューロンの遅い抑制性電位は、GABAの代謝型受容体の選択的阻害薬で消える。ドーパミン受容体には逆の作用をもつ二つのファミリーがある。線条体では、D1受容体をもつニューロンはシナプスの増強にドーパミンを必要とし、D2受容体をもつニューロンは減弱にドーパミンを必要とする & \cite{Colquhoun1992,DutarNicoll1988,Beaulieu2011,Shen2008} & 測定済み \\
12 & ブロードキャストチャネルは1本の配線ではなく、少数の線の束である（\ref{sec:buses}節） & マウス背側縫線核のセロトニンニューロンをウイルスと遺伝学的手法で標識：扁桃体に出力するニューロンと前頭皮質に出力するニューロンは核内の異なる場所にあり、互いに補い合う領域に枝分かれし、不快な刺激に逆向きに応答する。総説：青斑核（\nt{NORA}）のニューロンの亜群は異なる過程を符号化する & \cite{Ren2018,Poe2020} & 測定済み \\
13 & \nt{NORA}がゲイン$g$を決める & 適応的ゲインの仮説。カテコールアミンが伝達特性の傾きを上げるネットワークモデル。ネットワーク全体で「ノルアドレナリンとともに$g$が上がる」ことの直接測定はない。測定されているのは、サルの瞳孔径が青斑核ニューロンのスパイクに追従すること & \cite{AstonJones2005,ServanSchreiber1990,Joshi2016} & 仮説 \\
14 & \nt{ACET}が「書き込み/読み出し」を切り替え、学習率を決める & 仮説。根拠：ラット嗅皮質のスライスで、アセチルコリン作動薬（$100$\,\textmu M）は内部の「想起」シナプスを$60\,\%$以上弱めるが、入力シナプスは$15\,\%$未満しか弱めない & \cite{Hasselmo2006,HasselmoBower1992} & 仮説 \\
15 & \nt{SERO}が割引率$\gamma$を決める。セロトニンニューロンは一様に毎秒数スパイクで発火し、睡眠のある相ではほぼ沈黙する & 「セロトニン$=\gamma$」仮説。最も強い論拠：マウス背側縫線核のセロトニンニューロンを光遺伝学的に活性化すると、遅延報酬を待つのを途中でやめる回数が減る。発火率と急速眼球運動睡眠中の沈黙は測定されている（記録の総説） & \cite{Doya2002,Miyazaki2014,JacobsAzmitia1992} & 仮説 \\
\bottomrule
\end{longtable}}

\section{第\ref{sec:glutcomputer}章　\nt{GLUT}ニューロンは何を計算するか}

この章の論理的な主張は、非負の重みからなる回路についての定理であり、章の中で導いたものである。実験が裏づけるのは、ニューロンのモデルと、これらの定理に従って組んだ回路（同時検出器、遅延線、自己維持する群、連鎖）が脳に実際に存在することである。

{\footnotesize
\setlength{\tabcolsep}{3pt}\renewcommand{\arraystretch}{1.15}
\begin{longtable}{>{\raggedright}p{0.5cm}>{\raggedright}p{4.0cm}>{\raggedright}p{5.6cm}>{\raggedright}p{2.3cm}>{\raggedright\arraybackslash}p{1.7cm}}
\toprule
No. & 主張 & 実験と測定されたもの & 出典 & 状態 \\
\midrule\endfirsthead
\toprule
No. & 主張 & 実験と測定されたもの & 出典 & 状態 \\
\midrule\endhead
1 & ニューロンは閾値とリセットをもつRC回路である~\eqref{eq:glutneuron} & ラット皮質の錐体ニューロン（スライス）の細胞体に、生きた脳のシナプス入力の流れに似た雑音電流を注入した。平均発火率の電流の平均とばらつきへの依存性は、遅い発火率適応を加えれば漏れ積分発火ニューロンで記述できる。$\tau$と遅延の範囲は実験を引用せずに章で示した & \cite{Rauch2003} & 測定済み \\
2 & モデル~\eqref{eq:glutneuron}は単純化であり、入力の枝そのものが局所スパイクを出す & マウス視覚野の錐体ニューロンの樹状突起からin vivoで直接記録：視覚刺激がNMDA受容体に依存する局所的な樹状突起スパイクを起こし、それを抑えるとニューロンの方位選択性が広がる & \cite{Smith2013} & 測定済み \\
3 & 同時性の窓$\Delta_{\max}=\tau\ln\frac{w}{\theta-w}$~\eqref{eq:window}。$\theta=1.5w$では$0.7\tau$ & モデル~\eqref{eq:glutneuron}からの帰結。狭い加算窓の直接測定は速い抑制をもつニューロンについてある（第\ref{sec:gaba}章、その表の3行目） & 章内の導出 & 理論 \\
4 & \nt{GLUT}ニューロンだけのネットワークは入力の単調関数しか計算しない（命題~\ref{prop:monotone}） & 章内の証明：非減少の右辺による静寂からの反復。これはモデルについての主張であり、抑制のない生きたネットワークで検証した実験はない & 章内の導出 & 理論 \\
5 & 感覚器は配線ペアを与える。刺激の増加に応答する細胞と減少に応答する細胞がある & 網膜：双極細胞の段階ですでに錐体の信号はオンチャネルとオフチャネルに分かれる。サルでオン双極細胞を薬理学的に遮断：両チャネルは皮質まで別々に進み、遮断後は光の増加の検出が悪化するが、減少の検出は悪化しない & \cite{Schiller1992} & 測定済み \\
6 & 2線式符号があれば、\nt{GLUT}ニューロンのネットワークは任意の論理関数を共通のクロックに頼らず非同期に計算でき、代償はニューロン数が2倍になること（命題~\ref{prop:dualrail}、~\eqref{eq:dualrail}、図~\ref{fig:dualxor}） & 2線式符号と信号自体による完了検出は、遅延に依存しない非同期回路の標準的な技法である。ニューロン6個の排他的OR回路は章の中で組んだ。脳がまさに2線式符号で論理関数を計算していることを示す実験はない & \cite{Sparso2001} & 理論 \\
7 & ヒトで両耳への音の到着時刻の差は最大$\approx0.6\ms$であり、脳は約10マイクロ秒を区別する & 二肢選択実験の聴取者：時間差の弁別閾（正答率75\,\%）は、帯域$150$--$1700$\,Hzの雑音で$9$\,\textmu s、$1$\,kHzの純音で$11$\,\textmu s、クリック音で$28$\,\textmu s。上限$0.6\ms$は章の計算、$0.22\,\text{m}/343\,\text{m/s}$ & \cite{KlumppEady1956} & 測定済み \\
8 & 鳥では遅延線と同時検出器によって時間差が場所に変わる~\eqref{eq:itd}（図~\ref{fig:itd}） & メンフクロウ：両耳からの軸索は同時検出器の核に反対側から入る。軸索に沿ったスパイクの到着時刻は深さとともに規則的に変わり、核を横切る遅延の幅（$700$\,\textmu mで$160$\,\textmu s）は両耳間遅延の範囲と一致する & \cite{Carr1990} & 測定済み \\
9 & 哺乳類では遅延の列は見つかっていない。同じ課題を、\nt{GLYC}ニューロンからの時間的に精密な抑制をもつ同時検出器が解く & スナネズミの内側上オリーブでのin vivo記録：応答は方向の地図をなさない。マイクロピペットでグリシンとその阻害薬を投与すると、時間的に精密な\emph{グリシン}抑制が遅延の動作範囲を決めていることがわかる。ここでの抑制は\nt{GABA}ではなく\nt{GLYC}による & \cite{Brand2002,Grothe2010} & 測定済み \\
10 & 強い帰還をもつ群はフリップフロップであり、自己維持活動は刺激の後も数秒続く（図~\ref{fig:bistable}） & 記憶した点への遅延眼球運動課題（遅延$1$--$6$\,s）中のサル前頭前皮質：課題に関係する$170$個のニューロンのうち$87$個が遅延中に発火率を変え、そのうち$79\,\%$では記憶した点の方向に依存する。この活動を二つの安定状態をもつ帰還が維持しているというのは理論である & \cite{Funahashi1989,Wang2001} & 測定済み \\
11 & 流入が$\approx0.7\tau$より長く続けばビットが書き込まれる（図~\ref{fig:recurrentgroup}(b)） & 方程式$\tau\,\dd r/\dd t=-r+f(wr+I)$を$w=1$、$I=0.6$でオイラー法により計算。\texttt{models/}にスクリプトはない。実験はない & 章内の計算 & モデル \\
12 & 記憶はほとんどスパイクなしにシナプスの促通に保存できる & 理論：終末の残留カルシウムが書き込んだ内容を約1秒保持する。根拠：フェレットの内側前頭前皮質（複数ニューロンの同時記録）には、後効果の長い促通性シナプスで結ばれた錐体ニューロンのサブネットワークがある。記憶保持中の「静かな」区間の観察の総説。皮質がいつどちらの方法を使うかは未解決の問題 & \cite{Mongillo2008,WangMarkram2006,Stokes2015} & 仮説 \\
13 & 記憶は疲労によって消去される。神経伝達物質の蓄えが枯渇する & 皮質錐体ニューロンのペア記録：高頻度のスパイクでシナプスの応答が低下し、低下の速さは放出確率で決まる。これこそが自己維持する群を消すことは直接には示されていない & \cite{Tsodyks1997} & 測定済み \\
14 & 群の連鎖はイベントで起動するタイマーであり、鳴禽ではニューロンが厳密に順番に発火する & 睡眠中と歌唱中のキンカチョウの高次発声中枢で、同定されたニューロンを記録：運動核に出力を送る各ニューロンは、歌の一つの正確な瞬間に短いバーストを1回出し、ニューロンは順に発火する & \cite{Hahnloser2002} & 測定済み \\
\bottomrule
\end{longtable}}

\section{第\ref{sec:gaba}章　\nt{GLUT}と\nt{GABA}の協働}

この章の論理的・力学的な主張は、線形解析とオームの法則から章の中で導いたものである。実験が示すのは、それらが依拠する回路、つまり遅延つきの禁止、興奮に続く直接の抑制、抑制による安定化、\nt{GLUT}--\nt{GABA}ループのリズムが、脳で実際に働いていることである。

{\footnotesize
\setlength{\tabcolsep}{3pt}\renewcommand{\arraystretch}{1.15}
\begin{longtable}{>{\raggedright}p{0.5cm}>{\raggedright}p{4.0cm}>{\raggedright}p{5.6cm}>{\raggedright}p{2.3cm}>{\raggedright\arraybackslash}p{1.7cm}}
\toprule
No. & 主張 & 実験と測定されたもの & 出典 & 状態 \\
\midrule\endfirsthead
\toprule
No. & 主張 & 実験と測定されたもの & 出典 & 状態 \\
\midrule\endhead
1 & \nt{GABA}ニューロンは皮質のニューロンの5分の1から4分の1 & サル皮質の$10$領域でGABAを免疫染色：$8$領域でニューロンの約$25\,\%$、一次視覚野で$20\,\%$未満。皮質の抑制性ニューロンの多様性についての総説 & \cite{Hendry1987,Markram2004} & 測定済み \\
2 & 抑制が何をするかは着地点で大きく決まる。樹状突起、細胞体、スパイクの発生点、他の配線の終末 & 総説：皮質の\nt{GABA}ニューロンのクラスは受け手上の標的で区別される。錐体ニューロンの細胞体と樹状突起からの同時記録：直接の抑制は樹状突起より細胞体ではるかに強い。脊髄で他の配線の終末への抑制は、1本の入力線の神経伝達物質放出を減らす（測定の総説） & \cite{Markram2004,Pouille2001,Rudomin1999} & 測定済み \\
3 & 仲介役を通じた符号の反転には遅延1段、約1ミリ秒がかかる。興奮に続いて届く抑制は同時性の窓を狭める & ラット海馬スライスの錐体ニューロン：仲介役1個を介した抑制は直接の興奮に短い遅延で続き、二つの興奮性入力が閾値まで加算される窓は$2\ms$未満。この抑制が弱い樹状突起では窓が広い & \cite{Pouille2001} & 測定済み \\
4 & 禁止$a\wedge\bar b$~\eqref{eq:veto}はORと定数1とともに関数的に完全である（命題~\ref{prop:gabacomplete}） & 章内の証明。古典的な結果：抑制性入力をもつ閾値素子のネットワークは任意の論理式を実現する。モデルについての主張であり、実験はない & \cite{McCullochPitts1943} & 理論 \\
5 & 遅延つきの禁止は、最適速度$v^*=\Delta x/d$~\eqref{eq:vstar}の運動方向検出器になる & ウサギの網膜：方向選択性は、「ヌル」方向の運動で興奮を打ち消す抑制によって作られる。選択性をもつ細胞の興奮性入力と抑制性入力を別々に記録すると、スターバーストアマクリン細胞（\nt{GABA}）が非対称に、つまり「ヌル」側に向いた突起だけで抑制していることがわかった。$v^*$の式は章の導出 & \cite{Barlow1965,Fried2002} & 測定済み \\
6 & シャント抑制は電圧から引くのではなく割る~\eqref{eq:shunt}（図~\ref{fig:shunt}） & 塩化物イオンの平衡電位が静止電位付近にあるときの、三つのコンダクタンスからなる分圧器のオームの法則。スパイクを出さないニューロンについて & 章内の導出 & 理論 \\
7 & スパイクの発火率に対してシャントは引き算として作用し、割り算は抑制と釣り合った背景ノイズの組み合わせから得られる & モデル：発火しているニューロンではシャントを通る電流はほぼ一定で、発火率への効果は減算的。実験：ラット体性感覚野の錐体ニューロン（スライス）にダイナミッククランプで背景の興奮性・抑制性コンダクタンスを注入した。両者を釣り合わせて同時に増やすと、発火率をあまりずらさずに応答の傾きが割り算される。総説：全体の活動による割り算は脳の多くの部位で見られる & \cite{HoltKoch1997,Chance2002,Carandini2012} & 測定済み \\
8 & $\beta>1$では相互抑制がただ一つの勝者を選ぶ（命題~\ref{prop:wta}、~\eqref{eq:wta}） & 固有値$-(1\pm\beta)/\tau$は章の導出。$\beta=0.5$での発火率$0.73$と$0.53$は固定点$r_{1,2}=(I_{1,2}-\beta I_{2,1})/(1-\beta^2)$。\texttt{models/}にスクリプトはない。章では直接の実験は挙げていない & 章内の導出 & 理論 \\
9 & \nt{GABA}を介した信号によるメモリのリセット。相互抑制する二つの群はラッチである & 図~\ref{fig:bistable}と命題~\ref{prop:wta}から従う。実験はない & 章内の導出 & 理論 \\
10 & \nt{GLUT}--\nt{GABA}ループの平衡は、抑制が十分に強く十分に速ければ安定である（命題~\ref{prop:ei}） & 二集団モデル~\eqref{eq:ei}。条件(i)と(ii)はその行列の行列式とトレースであり、章で導いた & \cite{WilsonCowan1972} & 理論 \\
11 & $w_{EE}=1.5$、$w_{EI}w_{IE}=2$のとき、速い抑制（$\tau_I=2\ms$）は$E^*=0.67h$を保ち、遅い抑制（$30\ms$）は振動を起こす（図~\ref{fig:ei}） & 式~\eqref{eq:ei}の数値解、スクリプト\texttt{figures/make\_ei.py} & 計算 & モデル \\
12 & 皮質は抑制安定化の状態で動作する。その特徴は、\nt{GABA}ニューロンを後押しするとその発火率が下がること~\eqref{eq:isn} & 理論が逆説的な応答を予言した。実験：ネコ一次視覚野の細胞内記録で、周囲の刺激で応答が抑えられるとき、抑制性コンダクタンスは増えずに興奮性コンダクタンスと一緒に減る。マウスの視覚野、体性感覚野、運動野でPVニューロンを光遺伝学的に興奮させると、刺激がなくても、また軽い麻酔下でも、抑制性ニューロンの発火率が下がる。図~\ref{fig:ei}の値での係数$-1/3$は章の導出 & \cite{Tsodyks1997isn,Ozeki2009,Sanzeni2020} & 測定済み \\
13 & 「\nt{GLUT}が\nt{GABA}を興奮させ、\nt{GABA}が\nt{GLUT}を抑制する」ループは周期$12$--$50\ms$の速いリズムを生む & マウス体性感覚野の速い\nt{GABA}ニューロンを$8$--$200$\,Hzで光遺伝学的に刺激すると、ガンマリズム（$20$--$80$\,Hz、周期$12$--$50\ms$）が選択的に強まる。\nt{GLUT}ニューロンの刺激はより遅いリズムだけを強める & \cite{Cardin2009,Buzsaki2012} & 測定済み \\
\bottomrule
\end{longtable}}

\section{第\ref{sec:code}章　モールス符号}

この章の上限の見積もりは情報理論と計算によるものである。測定されているのは、通信路のどこでノイズが生じるか、脳がどれだけ速く動作するか、スパイクがどれだけのコストをもつかである。皮質が実際にどの符号を使っているかという問いは未解決のままである。

{\footnotesize
\setlength{\tabcolsep}{3pt}\renewcommand{\arraystretch}{1.15}
\begin{longtable}{>{\raggedright}p{0.5cm}>{\raggedright}p{4.0cm}>{\raggedright}p{5.6cm}>{\raggedright}p{2.3cm}>{\raggedright\arraybackslash}p{1.7cm}}
\toprule
No. & 主張 & 実験と測定されたもの & 出典 & 状態 \\
\midrule\endfirsthead
\toprule
No. & 主張 & 実験と測定されたもの & 出典 & 状態 \\
\midrule\endhead
1 & 皮質の\nt{GLUT}ニューロンの発火率は毎秒1スパイク未満から数十スパイクで、ニューロンはほとんどの時間を下限の近くで動作する & 覚醒ラットの聴覚野で、細胞に密着させたピペットによる記録（ニューロンの選択はその活動によらない）：各瞬間に毎秒$20$スパイクを超える発火率を示すニューロンは$5\,\%$未満。一部のニューロンの背景発火率は毎秒$0.01$未満。発火率の分布は対数正規 & \cite{Hromadka2008} & 測定済み \\
2 & スパイク通信路の容量は$R_{\max}\approx r\log_2\frac{e}{r\Delta t}$~\eqref{eq:spikeentropy}で、そのほとんどすべてがスパイクの時刻にある（命題~\ref{prop:capacity}） & 長さ$\Delta t$のセルからなる2進文字列のエントロピー。章の数値：$r=10$スパイク/s、$\Delta t=1\ms$でスパイク1個あたり$8$ビット、$\Delta t=10\ms$で約$5$ビット、スパイクカウンタでは毎秒$2$ビット未満 & \cite{MacKay1952} & 理論 \\
3 & 感覚ニューロンはスパイク1個あたり1ビットから数ビットを伝え、最もよい場合で上限の半分に達する & 刺激がわかっている感覚ニューロンのスパイクから刺激を復元。測定のまとめはこの書籍にある & \cite{Rieke1997} & 測定済み \\
4 & スパイク発生器は精確である。精確な時刻を決めるのは入力の速い変化である & ラット皮質のスライスのニューロン：速く揺らぐ電流を繰り返し与えるとスパイクは$1\ms$未満の精度で再現され、一定の電流では時刻がばらける & \cite{Mainen1995} & 測定済み \\
5 & シナプスは当てにならない。放出のランダム性のため、スパイクの半分以上は受け手に届かない & 海馬スライスで錐体ニューロンの入力を最小刺激：スパイクの半分以上が応答を起こさない。刺激閾値の揺らぎ、軸索分岐での伝導の失敗、実験温度の低さは除外された。残るのは確率的な放出である & \cite{AllenStevens1994} & 測定済み \\
6 & 生きた皮質のスパイク列はランダムに見える。間隔はポアソン列のようにばらつき、スパイク数の分散は平均に近い。「ノイズ」の一部は動物の運動についてのメッセージである & 覚醒サルの視覚野V1とMTでの記録：間隔の変動係数は約$1$、スパイク数の分散と平均の比はランダム過程と同様。独立な小さな入力を足し合わせるモデルでは、ばらつきははるかに小さくなる。マウス視覚野では、ばらつきのかなりの部分が動物自身の運動のビデオ記録から予測できる & \cite{Softky1993,Stringer2019} & 測定済み \\
7 & レート符号は遅い。誤差は$1/\sqrt{rT}$~\eqref{eq:rateerror}だが、脳は画像を$\sim150\ms$で認識する & 式はポアソン統計で、章の導出。実験：$20\ms$だけ見せた見知らぬ写真に動物がいるかどうかを被験者が判断した。「はい」と「いいえ」の脳電位は約$150\ms$後に分かれる。この時間を$10$--$20\ms$ずつ約10段に分けるのは章の推定であり、測定ではない & \cite{Thorpe1996} & 測定済み \\
8 & 群での平均化は相関によって制限される。誤差は$1/\sqrt c$分の1より小さくならない~\eqref{eq:correlated}（命題~\ref{prop:pooling}） & 同時記録したサルMT野のニューロンのペア：スパイク数の相関係数は平均$0.12$で、理論解析によればこの程度の相関でも群による改善は制限される。理論：限界を決めるのは相関のうち信号に似た部分だけである。反対の立場のモデル：$50$--$100$個のニューロンの群の発火率は1スパイク間隔、$10$--$50\ms$で読まれ、個々のスパイクの時刻は皮質では使われない & \cite{Zohary1994,MorenoBote2014,ShadlenNewsome1998} & 測定済み \\
9 & 数は最初のスパイクの時刻~\eqref{eq:latency}、群のスパイクの順序、局所リズムの位相に書くことができる & サンショウウオの網膜：短時間提示した画像の空間構造は、出力ニューロンの最初のスパイクの相対遅延に書かれており、コントラストにほとんどよらない。ラット海馬の場所細胞：「自分の」場所を通過するにつれて、スパイクはシータリズム（$7$--$12$\,Hz）のサイクルのますます早い時点へずれ、ずれは$100^\circ$から$355^\circ$に及ぶ。皮質がスパイクの時刻をどれだけ使うかは未解決の問題（8行目） & \cite{Gollisch2008,OKeefe1993} & 測定済み \\
10 & どの符号が読まれるかは受け手の時定数が決める~\eqref{eq:reader}（命題~\ref{prop:reader}）。活動中の皮質ではニューロンは同時検出器に近い & 式~\eqref{eq:reader}は章の導出。in vivo記録の総説：活動中の皮質では膜のコンダクタンスが静止時の数倍になり、時定数はそれに応じて短い。皮質がまさにこのように符号を切り替えることは直接には示されていない & \cite{Destexhe2003} & 仮説 \\
11 & 枯渇するシナプスは発火率の変化、本質的には$\Delta r/r$を伝える~\eqref{eq:depression}（図~\ref{fig:depression}） & 皮質錐体ニューロンのペア記録：枯渇の速さは放出確率で決まる。枯渇が遅いと応答は発火率を反映し、速いと同時性を反映する。ラット皮質で測定された枯渇に基づくモデル：速い入力と遅い入力で発火率の相対変化が等しければ応答も等しく、つまりシナプスは$\Delta r/r$を伝える。数値$1.7\to2.2$、$6.7$、$90\ms$は章の計算 & \cite{Tsodyks1997,Abbott1997depression} & 測定済み \\
12 & バーストは「ツー」である。失われる確率$(1-p)^k$~\eqref{eq:burst}は小さく、促通がそれをさらに下げる & 総説：当てにならないシナプスでは単発のスパイクはしばしば失われるが、バーストは促通のおかげで確実に伝わる。バーストはシナプスの長期的な変化を起こす。脳がバーストを独立した記号として読むというのは仮説 & \cite{Lisman1997,AllenStevens1994} & 仮説 \\
13 & 速い\nt{GABA}ニューロンはリズムと「句読点」を与える。もう一つのクラスは抑制するものを抑制し、ゲートを開く（命題~\ref{prop:punctuation}） & 皮質の\nt{GABA}ニューロンのクラスの性質の総説。光遺伝学：速い\nt{GABA}ニューロンのリズミックな興奮はガンマリズムを起こし、刺激への応答はサイクルの位相に依存する。マウス視覚野では、PVニューロンは互いを抑制し、SSTニューロンは他のすべてのクラスを、VIPニューロンはおもにSSTを抑制する。覚醒マウスでVIPニューロンを興奮させると\nt{GLUT}ニューロンの抑制が解除される。VIPニューロンをオンにするのは強化の信号である。一般規則としての役割分担は仮説 & \cite{Tremblay2016,Cardin2009,Pfeffer2013,Pi2013} & 仮説 \\
14 & エネルギーは平均発火率を毎秒1スパイク程度に制限する~\eqref{eq:energy} & 解剖学的・生理学的データによるげっ歯類灰白質のエネルギー収支：スパイクとシナプス電流は信号伝達のエネルギーの$47\,\%$と$34\,\%$。同時に活動できるニューロンは$\sim15\,\%$以下。ヒトの皮質では、同時に強く活動できるニューロンは$\sim1\,\%$以下。スパイク1個あたり$\sim10^9$個のATPと信号伝達の電力$\sim10$\,Wは、これらの計算に基づく章の推定 & \cite{Attwell2001,Lennie2003} & 理論 \\
15 & 符号はスパースで、1ジュールあたりのビット数が最大になるよう調整されている。皮質は精度をエネルギーと引き換えにする（命題~\ref{prop:economy}） & 経済的な符号の仮説。根拠：食餌制限したマウスの視覚野では、AMPA受容体のコンダクタンスとシナプスのATP消費が下がり（$-29\,\%$）、スパイクの発火率は保たれ、方位選択性は$32\,\%$広がる & \cite{Barlow1961,Padamsey2022} & 仮説 \\
\bottomrule
\end{longtable}}

\section{第\ref{sec:serocomputer}章　\nt{SERO}ニューロン}

\nt{SERO}ニューロン自体の性質（リズム、自己抑制、受容体による符号、サブシステム）は直接測定されている。本章の計算（調節器、フィルタ、論理）は本章の方程式から導かれる。$\tau_c$、セロトニンの拡散係数、放出部位の数は推定である。脳でループがレベル調節器として働くことは仮説である（縫線核では自己抑制は一対一で、短い休止を生むだけである）。\nt{SERO}が時間幅を担うことも仮説だが、これを支持する行動実験がある。

{\footnotesize
\setlength{\tabcolsep}{3pt}\renewcommand{\arraystretch}{1.15}
\begin{longtable}{>{\raggedright}p{0.5cm}>{\raggedright}p{4.0cm}>{\raggedright}p{5.6cm}>{\raggedright}p{2.3cm}>{\raggedright\arraybackslash}p{1.7cm}}
\toprule
No. & 主張 & 実験と測定内容 & 出典 & 状態 \\
\midrule\endfirsthead
\toprule
No. & 主張 & 実験と測定内容 & 出典 & 状態 \\
\midrule\endhead
1 & セロトニンの作用の符号は受け手の受容体が選ぶ（命題~\ref{prop:receptor}） & セロトニン受容体は薬理と機構によってファミリーに分けられる。5-HT$_{1A}$はGタンパク質を介してカリウムチャネルを開き、細胞を抑制する。5-HT$_{2A}$とイオンチャネル型の5-HT$_3$は興奮させる。一つの伝達物質が細胞ごとに逆の応答を生む。 & \cite{BarnesSharp1999} & 測定済み \\
2 & 覚醒時の\nt{SERO}ニューロンは毎秒数回のスパイクを一定に出す & 自由行動下のネコの背側縫線核ニューロンの記録。静かな覚醒時に$2.8\pm0.2$スパイク/sのゆっくりした規則的発火。徐波睡眠で発火率は$34$--$68\,\%$下がり、レム睡眠では$98\,\%$下がる。麻酔下のラットでは$1.7\pm0.5$スパイク/sで、非常に規則的。 & \cite{TrulsonJacobs1979, GagnonParent2014, JacobsAzmitia1992} & 測定済み \\
3 & \nt{SERO}ニューロンはペースメーカーである。スパイクごとの後押しは要らないが、一定の背景入力は要る（スライスでは$\alpha_1$受容体経由のノルアドレナリン、生体では\nt{NORA}から） & 脳スライスでは細胞はゆっくり規則的に発火し、各スパイクの後に$200$--$800\ms$の後過分極がある。スライスで自発発火する細胞は生体より少ない。規則的なリズムには$\alpha_1$受容体を介したノルアドレナリンの持続的入力が必要で、その遮断でリズムは消える。 & \cite{VandermaelenAghajanian1983} & 測定済み \\
4 & 受容体はほぼすべて代謝型。応答は遅く、およそ1秒のうちに立ち上がって減衰する & 5-HT$_3$を除くすべての受容体ファミリーはGタンパク質共役型である。スライスでは、誘発されたセロトニン放出が5-HT$_{1A}$を介して抑制電流を生み、それは1秒以内に立ち上がって減衰する。 & \cite{BarnesSharp1999, CourtneyFord2016} & 測定済み \\
5 & 投射先領域ではセロトニンは間隙だけでなく体積に出る。縫線核自体では近隣の抑制は一対一 & スライスでの高速サイクリックボルタンメトリー。1スパイクあたりの放出量は、シナプスのある領域と、シナプスがほとんどない縫線核とで同じであり、取り込みや受容体の遮断に依存しない。終末あたりの分子数はトランスポーターや受容体の数よりはるかに少なく、細胞外濃度は主要受容体の親和性と一致する。縫線核では5-HT$_{1A}$を介した抑制は一対一で、トランスポーターがセロトニンを間隙外に溜めない。 & \cite{Bunin1998, CourtneyFord2016} & 測定済み \\
6 & \eqref{eq:seroneuron}の濃度は回収によって減る。$\tau_c$が1秒程度というのは推定 & 生体脳でのボルタンメトリー。細胞外セロトニンを主に制限するのは合成ではなく、再取り込みと分解である。引用した研究に$\tau_c$の直接測定はなく、その大きさは本章の推定である。 & \cite{Hashemi2012serotonin} & 測定済み。$\tau_c$は推定 \\
7 & ペースメーカー1個でNOTとNOR。\nt{SERO}論理の完全性（命題~\ref{prop:serocomplete}、図~\ref{fig:serogates}） & \eqref{eq:seroneuron}から導かれる。抑制されるペースメーカーはNORであり、NORは関数的に完全である。\nt{SERO}ニューロンで論理を組んだ実験はない。体積が分かれているという条件は脳では満たされない。 & 本章で導出 & 理論 \\
8 & セロトニンは\nt{SERO}ニューロン自身をその上の受容体を介して抑制する & 麻酔下ラットで、5-HT$_{1A}$作動薬は静脈内投与でもイオントフォレシスでも、セロトニン自体と同じ用量反応関係で背側縫線核ニューロンの発火を抑制する。5-HT$_{1B}$作動薬はほとんど効かない。スライスでは、近隣細胞の内因性セロトニンが5-HT$_{1A}$を介して電流と発火の短い休止を生むが、平均発火率は変わらない。 & \cite{Sprouse1987, CourtneyFord2016} & 測定済み \\
9 & モデルでは、共通の体積を通るループはレベル調節器であり、ばらつきと時定数を$1+G$分の1にする~\eqref{eq:feedback}。脳でループがそう働くことは仮説 & 負帰還増幅器の古典的な式で、本章で導き直した。\nt{SERO}系のループゲイン$G$は測定されていない。測定されているのは、スライスでの自己抑制が一対一で、短い休止を生み、平均発火率を変えないことである。 & \cite{Black1934feedback, CourtneyFord2016} & 理論。脳では仮説 \\
10 & 濃度は発火率の移動平均、ローパスフィルタ~\eqref{eq:lowpass}、図~\ref{fig:lowpass} & 線形方程式~\eqref{eq:seroneuron}の解。図は$\tau_c=1\,\text{s}$でこの式から計算した。\texttt{models/}に個別のモデルはない。 & 本章で導出 & 理論 \\
11 & 間隙の屈曲は小分子の拡散をおよそ$2.6$倍遅くする & 点源法。スライスと生体脳でテトラメチルアンモニウムをイオントフォレシスで与え、その濃度をイオン選択性電極で記録する。屈曲度$\lambda\approx1.6$、つまり実効$D$は自由拡散の$\lambda^2\approx2.6$分の1。細胞外体積の割合は約$0.2$。セロトニン自体の組織内拡散係数はこれらの研究では測定されておらず、本章では推定である。 & \cite{Sykova2008} & 測定済み \\
12 & 独立な「配線」の長さ$\ell\sim\sqrt{D\tau}\approx20$~\textmu m~\eqref{eq:diffusionlength}。配線の数はそのような領域の数で制限される（命題~\ref{prop:volumewires}） & 拡散の法則に$D$と$\tau$の推定値（6行と11行）を代入したもの。命題~\ref{prop:volumewires}はその帰結である。容積伝達の独立チャネルを数えた直接の実験はない。 & 本章で導出 & 理論 \\
13 & 1個の\nt{SERO}ニューロンの出力は脳の広大な範囲に広がる。放出部位は推定$\sim10^5$ & ラット背側縫線核の単一ニューロンの完全再構成。上行軸索はすべて強く分岐し、1細胞の軸索の総長は最大$18.7$~cm、1細胞の枝が線条体、前頭前皮質、扁桃体に届く。細胞あたりの放出部位数は直接数えられていない。 & \cite{GagnonParent2014} & 測定済み。$10^5$は推定 \\
14 & \nt{SERO}ニューロンは出力と応答の異なるいくつかのサブシステムに分かれる & マウスでのウイルス遺伝学的標識。扁桃体に投射するニューロンと前頭皮質に投射するニューロンは分岐がほとんど重ならず、異なる入力を受け、不快刺激に逆向きに応答する。各群の活性化と抑制は行動を異なる形で変える。 & \cite{Ren2018} & 測定済み \\
15 & 指数割引での忍耐の条件$D<\tau_\gamma\ln(R/r_0)$~\eqref{eq:patience}。測定される双曲割引では$D<\tau_\gamma(R/r_0-1)$ & 割引$e^{-D/\tau_\gamma}$または$1/(1+D/\tau_\gamma)$から導かれる。導出は本章。秒単位の遅延での選択課題のサル。割引は指数より双曲線に近い。遅延報酬の予告信号に対する\nt{DOPA}ニューロンの応答も、遅延とともにやはり双曲線的に下がる。 & 本章で導出、\cite{KobayashiSchultz2008} & 理論 \\
16 & \nt{SERO}が時間幅$\tau_\gamma$を設定する（\ref{sec:horizon}節） & マウスで背側縫線核の\nt{SERO}ニューロンを光遺伝学的に活性化すると、遅延報酬を待つ時間が延び、まれになった報酬の探索の粘り強さが増す。ヒトではセロトニンの低下（トリプトファン欠乏食）で遅延報酬の割引が急になる。濃度がどのように$\tau_\gamma$に変わるかは分かっていない。 & \cite{Doya2002, Miyazaki2014, Lottem2018, Schweighofer2008} & 仮説 \\
\bottomrule
\end{longtable}}

\section{第\ref{sec:dofa}章　\nt{DOPA}による学習}

シナプスの機構（パケット、同時検出器、カルシウム、可塑性の窓）と、予測誤差としてのドーパミンのふるまいは直接測定されている。規則が勾配に導くこと、ブロードキャストの代償は定理である。選択学習の結果は本書のモデルによる計算である。誤差を計算する回路は仮説である。

{\footnotesize
\setlength{\tabcolsep}{3pt}\renewcommand{\arraystretch}{1.15}
\begin{longtable}{>{\raggedright}p{0.5cm}>{\raggedright}p{4.0cm}>{\raggedright}p{5.6cm}>{\raggedright}p{2.3cm}>{\raggedright\arraybackslash}p{1.7cm}}
\toprule
No. & 主張 & 実験と測定内容 & 出典 & 状態 \\
\midrule\endfirsthead
\toprule
No. & 主張 & 実験と測定内容 & 出典 & 状態 \\
\midrule\endhead
1 & 人工ネットワークにおける形の誤差逆伝播は脳では見つかっていない & 直接の実験はない。これは不在についての主張である。レビューは、このアルゴリズムが何を要求するか（順方向をなぞる逆向きの結合、独立したフェーズ）と、どんな生物学的近似が提案されているかを検討している。 & \cite{Lillicrap2020, WhittingtonBogacz2019} & 仮説 \\
2 & 重みは放出部位の数、放出確率、1パケットへの応答に分解される：$w=N_s\,p\,A_1$~\eqref{eq:quantal} & 放出を低下させたカエルの神経筋シナプス。受け手の応答は、1パケットに対する自発的な微小応答の整数倍の段階で変動する。スパイクあたりのパケット数はポアソン則に従う。 & \cite{delCastillo1954} & 測定済み \\
3 & 受け手の受容体は同時検出器であり、カルシウムが重みの変化を起動する & マウスニューロンのパッチクランプ。静止電位ではMg$^{2+}$イオンがグルタミン酸で開いたチャネルを塞ぎ、脱分極で外れる。海馬スライス。受け手内のカルシウムキレート剤は長期増強を阻止し、細胞内で光によりカルシウムを放出すると、それだけで伝達が増強される。 & \cite{Nowak1984, Malenka1988calcium} & 測定済み \\
4 & 決定はどちらの側でも実行される：受け手で$A_1$が増え、送り手で$p$が変わる & 1個のスパインの上で光により放出したグルタミン酸は、スパインとその受容体を通る電流の速く持続的な増大を起こす。増強にはAMPA受容体の挿入が伴う。受け手の脱分極は送り手の放出を一時的に抑える。その効果は間隙を逆向きに進むエンドカンナビノイドが運ぶ（\nt{GABA}入力で示された）。短期的な枯渇は送り手の放出確率に依存する。 & \cite{Matsuzaki2004, Malinow2002, WilsonNicoll2001, Tsodyks1997} & 測定済み \\
5 & 送り手と受け手が同時に活動するとシナプスは強まる~\eqref{eq:hebb} & 麻酔下のウサギ。海馬の入力路への短いスパイク列の後、応答は$30$~分から$10$~時間増強される。海馬スライスでは、受け手のスパイクは同じ瞬間に活動しているシナプスだけを強める。 & \cite{Bliss1973LTP, Kelso1986hebb} & 測定済み \\
6 & 規則~\eqref{eq:oja}は入力の相関の主固有ベクトルを見つける（命題~\ref{prop:oja}） & 定理。証明は本章。ニューロンが入力の主成分を学んだという実験はない。シナプスで重みの増大を制限する機構は確立していない。 & \cite{Oja1982} & 理論 \\
7 & タイミングによるヘッブ則：窓は約$\pm20\ms$（図~\ref{fig:stdp}）。繰り返される系列から連鎖が育つ & 培養した海馬ニューロンのペア。送り手のスパイクの後$20\ms$以内の受け手のスパイクは持続的な増強を、前$20\ms$以内のものは抑圧を生み、どちらもNMDA受容体に依存する。図の比$A_-=0.6A_+$は説明用である。ネットワークでこうして連鎖が育つことは直接測定されていない。 & \cite{BiPoo1998} & 測定済み \\
8 & トレース~\eqref{eq:threefactor}：同時の活動の後$1$--$2\,\text{s}$に届いたドーパミンは結合を強め、それより遅いものは強めない（命題~\ref{prop:threefactor}） & 線条体スライスで、グルタミン酸入力とドーパミン入力を別々に光遺伝学的に刺激。ドーパミンがグルタミン酸入力の$0.3$--$2\,\text{s}$後に来たときだけスパインが大きくなる。窓はcAMPの速い上昇と分解で決まる。皮質では、ペアリングが増強と抑圧の別々のトレースを残し、モノアミン受容体がそれを秒程度の窓で長期的な変化に変える。 & \cite{Yagishita2014, He2015eligibility} & 測定済み \\
9 & ドーパミンなしでも秒単位の窓がある：第三の信号は受け手の強い興奮 & 生きたマウスのCA1野。受け手の1回のプラトー事象が、その前後数秒に来た入力を強め、1回の通過で場所野を作る。 & \cite{Bittner2017} & 測定済み \\
10 & 三要素則の平均ステップは報酬の勾配に沿う~\eqref{eq:perturbation}（命題~\ref{prop:gradient}） & ランダムな摂動による学習の勾配定理。本章では小さな雑音について導いた。 & \cite{Williams1992} & 理論 \\
11 & 雑音は探索である：ネットワークは自分のランダムなずれから学ぶ & 若いキンカチョウ。大脳基底核回路から出る核を不活性化すると、歌の変動性が急激かつ可逆的に消える。成鳥のキンカチョウ。平均より一方の側の高さで歌われた音節に妨害を与えると、鳥はすばやく高さを反対側にずらす。自分のばらつきから学ぶ。 & \cite{Olveczky2005, TumerBrainard2007} & 測定済み \\
12 & 二者択一は$\tau_e=1\,\text{s}$、$\delta=R-\bar R$で学習され、$\tau_e=50\ms$では学習されない。期待値を引かないと重みは際限なく増え、学習率が大きいとネットワークは悪い選択を固定しうる & \texttt{models/dofa\_learning.py}、$\eta=0.05$、$500$試行、シード$6$個。$\tau_e=1\,\text{s}$：$A$を選ぶ割合$0.65$--$0.79\to0.96$--$1.00$、重み$0.85$--$1.33$。$\tau_e=50\ms$：$0.38$--$0.55$、重みは変わらない。$\delta=R$：勝者の重み$860$--$1190$。$\delta=R$、$\eta=0.5$、シード$3$個：1個が$B$に固定（$0.04\to0$）。スクリプトは四つの場合をすべて出力する。再計算は本章と一致した。 & \texttt{models/\allowbreak dofa\_\allowbreak learning.py} & モデル \\
13 & 共通の信号一つによる学習時間は、雑音を出すニューロンの数に比例して伸びる（命題~\ref{prop:broadcastcost}） & 線形ネットワークの学習曲線。ニューロン摂動は正確な勾配より出力ニューロン数倍遅く、重み摂動はさらに入力数倍遅い。 & \cite{Werfel2005} & 理論 \\
14 & \nt{DOPA}の出力は主に行動選択の領域に向かう & ラットの黒質線条体\nt{DOPA}ニューロン個々の軸索の完全再構成。1細胞の枝は線条体体積の$0.45$--$5.7\,\%$（平均$2.7\,\%$）を覆い、線条体の外にはほとんど枝がない。 & \cite{Matsuda2009} & 測定済み \\
15 & 皮質は自分の入力を予測することを学び、誤差はその場で計算される & レビュー。感覚皮質には予期した入力と実際の入力の不一致に対する応答がある。これが皮質の一般的な学習則であることは証明されていない。 & \cite{Keller2018} & 仮説 \\
16 & ドーパミンは誤差~\eqref{eq:ctd}のようにふるまう：バースト、予告信号への移行、落ち込み（図~\ref{fig:ctdcases}） & サルの\nt{DOPA}ニューロン。予期しないジュースにバースト。学習後は予告信号にバーストが出て、約束のジュースには応答しない。ジュースが来ないと落ち込む。連続形~\eqref{eq:ctd}は理論である。 & \cite{Schultz1997, Doya2000} & 測定済み \\
17 & $\dd V/\dd t$を計算し期待を差し引く回路 & マウスの腹側被蓋野での記録と光遺伝学。\nt{DOPA}ニューロンは報酬から期待を差し引く。近くの\nt{GABA}ニューロンは報酬が期待されるときにそれを抑制し、その興奮や抑制は誤差を変える。微分がどう計算されるかは示されていない。 & \cite{Eshel2015arithmetic} & 仮説 \\
18 & \nt{DOPA}ニューロンはペースメーカーで、落ち込みはバーストより浅い & スライスで\nt{DOPA}ニューロンは、ゆっくりしたペースメーカー電位の上で自発的に非常に規則的に発火する。サルではレートは正の誤差を定量的に追うが、負の誤差は弱くしか伝えない。 & \cite{GraceOnn1989, Bayer2005} & 測定済み \\
19 & アクターとクリティック（図~\ref{fig:actorcritic}） & アルゴリズムは強化学習の理論から。ヒト（fMRI）では、腹側線条体は選択があってもなくてもクリティックのようにふるまい、背側線条体は行動を選ぶ必要があるときだけそうふるまう。 & \cite{SuttonBarto2018, ODoherty2004actor} & 仮説 \\
20 & 第二の受容体ファミリーをもつニューロンでは規則の$\delta$の符号が反転する & 線条体スライス。D1受容体をもつニューロンではペアリングがD1を要する増強を、D2をもつニューロンではD2を要する抑圧を生む。 & \cite{Shen2008} & 測定済み \\
\bottomrule
\end{longtable}}

\section{第\ref{sec:dofaalt}章　\nt{DOPA}の他の理論}

拡張された描像はそれぞれ、ドーパミンの独自の直接測定（反転点のばらつき、不快なものと新しいものへの応答、ゆっくりした増加、味の変化への応答）に基づいている。しかしそれらを説明する式は理論であり、そのうち二つの間では実験がまだ互いに矛盾している。

{\footnotesize
\setlength{\tabcolsep}{3pt}\renewcommand{\arraystretch}{1.15}
\begin{longtable}{>{\raggedright}p{0.5cm}>{\raggedright}p{4.0cm}>{\raggedright}p{5.6cm}>{\raggedright}p{2.3cm}>{\raggedright\arraybackslash}p{1.7cm}}
\toprule
No. & 主張 & 実験と測定内容 & 出典 & 状態 \\
\midrule\endfirsthead
\toprule
No. & 主張 & 実験と測定内容 & 出典 & 状態 \\
\midrule\endhead
1 & 非対称な規則~\eqref{eq:asymmetric}は点~\eqref{eq:expectile}に収束する。$\kappa=1/2$ではそれは平均、確率$1/2$の1滴では$V=\kappa$（図~\ref{fig:expectiles}） & 点~\eqref{eq:expectile}はエクスペクタイル、つまり正と負の偏差に異なる重みをつけた最小二乗問題の解である。本章の例の導出は本章内にある。 & \cite{NeweyPowell1987}、本章で導出 & 理論 \\
2 & \nt{DOPA}ニューロンごとに反転点が異なり、非対称性の順に並ぶ（命題~\ref{prop:expectile}） & マウス。光遺伝学的に標識した腹側被蓋野の\nt{DOPA}ニューロン、大きさの異なる報酬。バーストが落ち込みに変わる点はニューロンごとに異なる。応答の非対称性が大きいニューロンほどその点は高い。群の応答から報酬分布の形が復元される。これがばらつきの主な機能であることは仮説。 & \cite{Dabney2020} & 測定済み \\
3 & \nt{DOPA}ニューロンの一部は不快なものにバーストで応答する & サル。ある\nt{DOPA}ニューロンは報酬の予告信号で興奮し、目への空気噴射の予告信号で抑制される。別のものは両方で興奮する。二つの群は中脳の異なる部分にある。 & \cite{Matsumoto2009, BrombergMartin2010} & 測定済み \\
4 & 誤差の絶対値または新奇性が学習率を決める & 引用した研究には、\nt{DOPA}ニューロンの重要性信号が学習率を変えることを示す直接の実験はない。レビューは「注意、重要」という信号の役割を提案している。 & \cite{BrombergMartin2010} & 仮説 \\
5 & 危険の軸のための別の配線 & マウス。線条体の後端に向かう\nt{DOPA}ニューロンは報酬ではなく新しい強い刺激に応答する。その破壊は新しい物体の回避を弱める。 & \cite{Menegas2018} & 測定済み \\
6 & 最適な行動時間$\tau_a^*=\sqrt{C/\bar R}$~\eqref{eq:vigor} & 和$C/\tau_a+\bar R\tau_a$の最小値。導出は本章。元のモデルでは、行動の速さは平均報酬に等しい時間のコストで決まる。 & \cite{Niv2007}、本章で導出 & 理論 \\
7 & ドーパミンのゆっくりしたレベルは$\bar R$を運び、行動の速さを決める（命題~\ref{prop:multiplex}） & ラット、側坐核でのマイクロダイアリシスとボルタンメトリー。ドーパミンのレベルは分単位で報酬の頻度と行動の速さに従う。ヒトではL-DOPAが反応時間の平均報酬頻度への依存を強める。遅いレベルがまさに$\bar R$に等しいことは証明されていない。 & \cite{Hamid2016, Beierholm2013vigor} & 仮説 \\
8 & ゆっくりしたレベルは二つの源からなる。放出はスパイクのレートを変えずに出力で変わるが、ゆっくりした増加はレート自体にもある & ラット、一つの課題。側坐核での放出は変化する報酬の期待に従うが、同定された腹側被蓋野\nt{DOPA}ニューロンのスパイクのレートは従わない。仮想通路のマウス（10行）。報酬への接近に伴うなめらかな増加は、同定された\nt{DOPA}ニューロンのスパイクにもある。 & \cite{Mohebi2019, Kim2020} & 測定済み \\
9 & 動物が遠くの報酬に近づく間、ドーパミンはなめらかに増える & 迷路のラット、線条体でのボルタンメトリー。ドーパミン濃度は目標に近づくにつれ数秒かけて上がり、その傾きは距離と報酬の大きさに依存する。 & \cite{Howe2013ramp, Hamid2016} & 測定済み \\
10 & 増加は推定の精密化に伴う誤差$\delta$である~\eqref{eq:ramp}。通路での跳躍はバーストを生む（図~\ref{fig:ramp}） & 式と毎秒$\delta=0.75\,V$という数値は本章で導出。実験：仮想通路のマウスを目標近くへ瞬時に移す。細胞体のスパイク、細胞体と出力のカルシウム、線条体での濃度は、期待$V$ではなく誤差のように応答する。二つの説明のどちらがすべての出力で正しいかは議論中。 & \cite{Kim2020} & 測定済み \\
11 & 振り返り：ドーパミンは因果的結びつきの尺度~\eqref{eq:retro}を伝える & マウス、この理論と誤差~\eqref{eq:ctd}が異なる予測をする実験での側坐核のドーパミン放出。いくつかの実験で放出は前者に従った。 & \cite{Jeong2022} & 仮説 \\
12 & 学習中のドーパミン応答は報酬から予告信号へ徐々に移る & マウス、学習の進行に伴う腹側線条体での\nt{DOPA}ニューロン出力の信号。バーストは\eqref{eq:ctd}による学習が要求するとおり、報酬の時点から予告信号の時点へ徐々に移る。 & \cite{Amo2022} & 測定済み \\
13 & \nt{DOPA}ニューロンは同じ価値での報酬の味の変化に応答する & ラット、腹側被蓋野ニューロンの記録。ジュースを価値の等しい別のものに替えると、誤差~\eqref{eq:ctd}はゼロなのに応答が生じる。 & \cite{Takahashi2017} & 測定済み \\
14 & 人工的なバーストは二つの中立な刺激の結びつきを教える & ラット、報酬なしに二つの刺激をあらかじめ対にする実験。第一から第二への移行時に\nt{DOPA}ニューロンを光遺伝学的に活性化すると結びつきが生じ、抑制するとその学習が妨げられる。 & \cite{Sharpe2017} & 測定済み \\
15 & 特徴ごとの誤差のベクトル~\eqref{eq:vectortd} & 特徴の予測誤差の理論。ベクトル形式の直接の検証はない。 & \cite{Gardner2018} & 仮説 \\
16 & 一つの課題の中で\nt{DOPA}ニューロンごとに異なる量を運ぶ & 仮想通路のマウス、\nt{DOPA}ニューロンの2光子イメージング。報酬に関係するもの、速度と旋回に関係するもの、予告信号と選択の正確さに関係するものがある。 & \cite{Engelhard2019} & 測定済み \\
17 & 配線ではなく配線の束（命題~\ref{prop:bundle}） & 2--16行のまとめ。各分解は個別に測定されている。それらがすべて一つの信号の部分であるという統一的な描像は実験で示されていない。 & --- & 仮説 \\
\bottomrule
\end{longtable}}

\section{第\ref{sec:nora}章　\nt{NORA}}

\nt{NORA}系の構成、その二つのモード、瞳孔との関係（\nt{NORA}経由だけではない）、信号対背景比の上昇、行動における逆U字は直接測定されている。乗算的なゲイン$g$はモデルである。ゲインが選択のモードを切り替え、逆温度として働くことは本章の導出である。ゲイン調整の利点は本書のモデルによる計算である。「\nt{NORA}は探索のつまみ」と「\nt{NORA}は割り込み」は仮説である。

{\footnotesize
\setlength{\tabcolsep}{3pt}\renewcommand{\arraystretch}{1.15}
\begin{longtable}{>{\raggedright}p{0.5cm}>{\raggedright}p{4.0cm}>{\raggedright}p{5.6cm}>{\raggedright}p{2.3cm}>{\raggedright\arraybackslash}p{1.7cm}}
\toprule
No. & 主張 & 実験と測定内容 & 出典 & 状態 \\
\midrule\endfirsthead
\toprule
No. & 主張 & 実験と測定内容 & 出典 & 状態 \\
\midrule\endhead
1 & ヒトの\nt{NORA}ニューロンは数万個で、ほぼすべてが一つの集団にある & ヒト脳幹の連続切片、チロシン水酸化酵素染色。青斑核に$53\,900$個、隣接する亜核にさらに$6\,260$個。 & \cite{Baker1989LC} & 測定済み \\
2 & 出力は脳のほぼ全体に広がるが、集団の部分ごとにある程度異なる領域を担当する & マウスでのウイルス遺伝学的標識。青斑核の\nt{NORA}ニューロンは脳の広い領域から入力を受け、脳と脊髄の広い領域に出力を送る。ラットでは、扁桃体に向かう群は恐怖の学習を強め、前頭前皮質に向かう別の群はその消去を強める。 & \cite{Schwarz2015LC, Uematsu2017LC, Poe2020} & 測定済み \\
3 & 持続モード：毎秒1回未満から数回の規則的なスパイクで、発火率は覚醒度とともに変わる & ラット、睡眠覚醒サイクルでの青斑核ニューロンの記録。発火率は覚醒時に最大、徐波睡眠で低く、レム睡眠でほぼ0。発火率の変化は状態の切り替わりに先行する。 & \cite{AstonJonesBloom1981} & 測定済み \\
4 & 相動モード：課題にとって重要な出来事に、ほぼ決定の瞬間に短いバースト & まれな標的を見つける課題のサル。青斑核のすべてのニューロンが標的にだけバーストで応答し、その$\sim90\ms$後、手の応答の$\sim200\ms$前に出る。成績が悪いときはバーストが弱い。二者択一の課題では、バーストは刺激より応答に正確に結びつき、正答でも誤答でも出る。 & \cite{AstonJones1994vigilance, Clayton2004LC} & 測定済み \\
5 & 瞳孔径は\nt{NORA}の不正確なテストポイント：\nt{NORA}にも\nt{ACET}にも他の領域にも従う & サル。青斑核のスパイクは、1細胞の個々のスパイクに至るまで瞳孔の散大を予測する。似ているがより遅い関係が丘と帯状皮質にもある。マウス。瞳孔のゆらぎは、皮質のノルアドレナリン作動性出力とコリン作動性出力の両方の活動に従う。 & \cite{Joshi2016, Reimer2016} & 測定済み \\
6 & ノルアドレナリンは背景活動を誘発活動より強く抑える。乗算的ゲイン~\eqref{eq:gain}はこの効果を再現するモデル & 覚醒サルの聴覚皮質、ノルアドレナリンのイオントフォレシス。ニューロンの背景活動は同種の声への応答より強く抑えられ、SN比が上がる。乗算的シグモイド~\eqref{eq:gain}はネットワークモデルからとったもので、傾きの文字どおりの乗算は測定されていない。 & \cite{Foote1975NE, ServanSchreiber1990} & 測定済み。$g$はモデル \\
7 & ゲインが$d'$を上げるのは増幅器の後のノイズに対してだけ~\eqref{eq:gainsnr}（命題~\ref{prop:gainnoise}） & 本章で導出。ネットワークモデルではゲインは信号検出を改善する。増幅器の前と後のノイズを分けてこの違いを検証した実験はない。 & 本章で導出、\cite{ServanSchreiber1990} & 理論 \\
8 & 境界$g\beta=1$は選択ネットワークを「思案中」から「決定済み」に切り替える~\eqref{eq:wtagain}（命題~\ref{prop:gainwta}、図~\ref{fig:pitchfork}） & 互いに抑制し合う二つの群の線形解析。本章で導出。脳でのこの閾値の直接測定はない。 & 本章で導出 & 理論 \\
9 & ゲインは選択の逆温度~\eqref{eq:softmax} & 増幅器の後のロジスティックノイズでは、選択確率は$T=\sigma/g$のソフトマックスになる。本章で導出。 & 本章で導出 & 理論 \\
10 & \nt{NORA}はどれだけ探索するかを決める：高い持続性活動は小さな実効$g$（探索）、決定の瞬間の相動性バーストは大きな実効$g$（決定） & ヒト。当て推量の選択の前には、既知の最良の選択の前よりベースラインの瞳孔が広い。瞳孔が広い人ほど探索の傾向が強い。ラット。ランダムな選択のモードへの移行は、前帯状皮質への\nt{NORA}入力が制御する。バーストが実効$g$を上げることは直接測定されていない。 & \cite{JepmaNieuwenhuis2011, Tervo2014stochastic, AstonJones2005} & 仮説 \\
11 & 報酬は$g$に対して逆U字に依存する。速く変わる世界では最良の$g$は小さい（命題~\ref{prop:explore}、図~\ref{fig:invertedU}） & \texttt{models/nora\_explore.py}、$4000$試行の実行$60$回。$1000$試行に1回の変化：最良$g=16$、割合$0.976$。$20$試行に1回：$g=8$、割合$0.886$。$g=0.5$では$0.77$。再計算は本章と一致した。 & \texttt{models/\allowbreak nora\_\allowbreak explore.py} & モデル \\
12 & 行動における逆U字：明瞭なバーストを伴う中程度の持続性活動で成績が最もよい & 4行と同じ課題のサル。成績のよい時期には持続性発火率が中程度で、標的へのバーストが明瞭である。持続性発火率が高いとバーストがなく、誤反応が多い。持続性活動と誘発活動の変化は成績のゆらぎに従う。 & \cite{Usher1999LC, AstonJones2005} & 測定済み \\
13 & \nt{NORA}は予期しない不確実性を、\nt{ACET}は予期された不確実性を伝える & 二種類の不確実性をもつベイズ推論の理論。引用した研究には、これらの信号を伝達物質で分ける直接の実験はない。 & \cite{YuDayan2005} & 仮説 \\
14 & 急な変化の後、人はより速く学び、瞳孔が広がる & ヒト、急な変化のある予測課題。瞳孔の短い変化は変化確率の推定に従い、ベースラインの瞳孔とともに、新しいデータが推定をどれだけ変えるか（学習率）を予測する。光や課題と無関係な瞳孔の変化もこの率を変える。ここで瞳孔がまさに\nt{NORA}を示すというのは、5行の間接的データからの推論である。 & \cite{Nassar2012} & 測定済み \\
15 & 相動性バーストは割り込みとして働く：ネットワークは状態をリセットする & \nt{NORA}のバーストがネットワークの状態をリセットする直接の実験はない。測定されているのは重要な出来事へのバーストだけである（4行）。 & \cite{BouretSara2005} & 仮説 \\
16 & 驚きによる$g$や学習率の調整は、最良の一定設定よりほとんどよくない & \texttt{models/nora\_explore.py}、各時期$1000$試行の世界（$1000$試行に1回と$20$試行に1回の変化）。最良の一定$g=8$：$0.925$。$g$の調整：最大$0.927$。学習率の調整：最大$0.926$。一定の学習率$0.4$：$0.928$。再計算は本章と一致した。 & \texttt{models/\allowbreak nora\_\allowbreak explore.py} & モデル \\
\bottomrule
\end{longtable}}

\section{第\ref{sec:acet}章　\nt{ACET}を加える}

アセチルコリンの三つの作用はスライスと生体脳で測定されている。書き込みと読み出しの衝突は本書のモデルで示した。\nt{ACET}が書き込み許可線として働き、予期された不確実性を伝えることは、部分的に実験の裏づけがある仮説である。

{\footnotesize
\setlength{\tabcolsep}{3pt}\renewcommand{\arraystretch}{1.15}
\begin{longtable}{>{\raggedright}p{0.5cm}>{\raggedright}p{4.0cm}>{\raggedright}p{5.6cm}>{\raggedright}p{2.3cm}>{\raggedright\arraybackslash}p{1.7cm}}
\toprule
No. & 主張 & 実験と測定内容 & 出典 & 状態 \\
\midrule\endfirsthead
\toprule
No. & 主張 & 実験と測定内容 & 出典 & 状態 \\
\midrule\endhead
1 & 脳の\nt{ACET}ニューロンは少なく、いくつかの群にまとまり、出力は皮質全体に広がる：ブロードキャストチャネル & ラットでのアセチルコリン合成酵素に対する抗体染色と逆行性標識。皮質、海馬、扁桃体、嗅球、視床は、主として前脳基底部と脳幹にある六つのコンパクトな細胞群からアセチルコリンを受け取る & \cite{Mesulam1983} & 測定済み \\
2 & 皮質のアセチルコリンレベルは覚醒時に高く、徐波睡眠で低い & 自由行動下のネコの皮質と海馬でのマイクロダイアリシスとEEG記録。アセチルコリン放出は徐波睡眠に比べ、静かな覚醒で$\sim100\%$、活動的な覚醒で$\sim175\%$高い。レム睡眠の皮質では覚醒時と同程度 & \cite{Marrosu1995,Hasselmo1999} & 測定済み \\
3 & 信号の意味は受容体が決める：アセチルコリンには速いチャネル型受容体と、細胞内の反応を起動する遅い受容体がある（命題~\ref{prop:receptor}） & 測定のレビュー。興奮性、伝達、可塑性に対するアセチルコリンの作用は、放出部位、受容体サブタイプ（ニコチン性はイオンチャネル、ムスカリン性はGタンパク質経由）、受け手の細胞に依存する & \cite{Picciotto2012} & 測定済み \\
4 & 第一の作用：外部入力にほとんど触れずに内部結合を弱める（\eqref{eq:acetnet}の係数$1-a$） & ラット嗅皮質のスライス。$100$\,\textmu Mのアセチルコリン作動薬で、同じ細胞において、内部結合（Ib層）の電位は$60\%$以上、外部入力（Ia層）の電位は$15\%$未満低下する。抑制はシナプス前性。海馬スライス（CA1）では内部経路と入力経路で$89\%$対$40\%$ & \cite{HasselmoBower1992,HasselmoSchnell1994} & 測定済み \\
5 & 第二の作用：入力を強める（係数$\frac{1+a}{2}$） & 新皮質スライス。ニコチン性受容体は視床からの入力のシナプスを強め、皮質内のシナプスには作用しない。ムスカリン性受容体は両方を弱める。アセチルコリンはニューロンの興奮性を高める（レビュー） & \cite{Gil1997,Hasselmo2006} & 測定済み \\
6 & 第三の作用：重みの変化を容易にする（学習率$\eta a$） & ラット。基底核（皮質へのアセチルコリンの供給源）の電気刺激を音と対にすると、成体でも聴覚皮質の地図が提示音に合わせて大きく再編される & \cite{KilgardMerzenich1998,Hasselmo2006} & 測定済み \\
7 & \eqref{eq:acetnet}の第二式のまばらなパターン用ヘッブ則は、一部からパターンを補完する連想記憶を与える & まばらなパターンと共分散則をもつネットワークの容量計算。活動ニューロンの割合$p$が小さいと、ネットワークは$p=1/2$の場合よりはるかに多くのパターンを保存する & \cite{Tsodyks1988} & 理論 \\
8 & 低い$a$での書き込みは記憶を壊す：ネットワークは見たものではなく思い出したものを書き込む。$a=0.1$での破損は$a=0$より弱くない & \texttt{models/\allowbreak acet\_memory.py}：ニューロン$200$個、$20$個ずつのパターン五つ、新パターンはそのうち一つと$10$個のニューロンを共有、$10$回の実行。$a=0$では新パターンは記憶されず、古いパターンは$10$個中$5.9$個の無関係なニューロンを引き連れる。$a=0.1$では新パターンは思い出される（$0.97$）が両者が融合し、新パターンは$7.3$個、古いパターンは$7.9$個の無関係なニューロンを引き連れる。$a=0.2$からは1個未満 & 本章で導出 & モデル \\
9 & 命題~\ref{prop:writeread}：書き込みと読み出しが同じようにうまくいくレベル$a$はない（図~\ref{fig:acetsim}） & 同じスクリプト、学習率$\eta a$。新パターンが1回の提示で丸ごと記憶されるのは$a\ge0.9$、$a=0.75$で$0.8$、$a\le0.5$では記憶されない。パターンの半分からの読み出し：$a\le0.2$で$1.0$、$a=0.3$で$0.96$、$a=0.4$で$0.04$ & 本章で導出 & モデル \\
10 & 脳では1本のブロードキャスト配線、\nt{ACET}が書き込みと読み出しを切り替える & この考えはスライスでの測定に基づくモデルからとった。最も直接的な実験はヒトでのもの。ムスカリン受容体遮断薬スコポラミンは新しい単語対の暗記、特に古いものと重なる対の暗記を悪化させるが、すでに学習した対の想起は悪化させない。ネットワークの二つのモードの直接測定はない & \cite{HasselmoAndersonBower1992,Atri2004} & 仮説 \\
11 & フィルタ~\eqref{eq:kalman}は最適である。入力の重みと学習率は同じ量$P/R$で、定常状態では$P=\sqrt{qR}$、記憶は$\sqrt{R/q}$（命題~\ref{prop:kalman}） & 実験を要しない。カルマン--ビューシーフィルタの最適性は推定理論の古典的結果であり、定常状態は本章で導いた & \cite{KalmanBucy1961} & 理論 \\
12 & \nt{ACET}はネットワークに$P$に似た量、予期された不確実性を伝える & 注意の確率モデルで提案された。アセチルコリンのレベルが推定の不確実性に従うことの直接測定はない & \cite{YuDayan2005} & 仮説 \\
13 & \nt{ACET}ニューロンの一部は報酬と罰に20ミリ秒ほどで応答し、強化が予期しないものであるほど強く応答する & マウス、光遺伝学的に同定した前脳基底部のコリン作動性ニューロン。報酬と罰への応答は$18\pm3$\,msで、大きさは強化の予期しなさとともに増し、二つの核のニューロンで応答は似ている & \cite{Hangya2015} & 測定済み \\
14 & \nt{NORA}は世界の予期しない変化に気づき、\nt{ACET}はネットワークを書き込みモードに保つ & この分業はモデルで提案された。間接的には、ヒトで\nt{NORA}に関係する瞳孔径が変化と学習率を追跡する。\nt{NORA}--\nt{ACET}の組の直接の検証はない & \cite{YuDayan2005,Nassar2012} & 仮説 \\
15 & 徐波睡眠では、読み出しモードのネットワークが日中に書き込んだものを再生し、その再生が長期記憶への書き直しになる & ラット海馬のニューロン集団の記録。日中に一緒に働いた細胞は、その後の徐波睡眠でより頻繁に、同じ順序で一緒に発火する（測定済み）。書き直しについて。学習後に海馬のリップルを抑制すると空間記憶が悪化し、自発リップルを延ばすと改善する。原因が\nt{ACET}の低レベルであることは証明されていない & \cite{WilsonMcNaughton1994,Skaggs1996,Girardeau2009,FernandezRuiz2019,Hasselmo1999} & 仮説 \\
\bottomrule
\end{longtable}}

\section{第\ref{sec:synthesis}章　マシン全体}

まとめの章なので新しい実験は挙げない。各行は、その主張を詳しく扱った章と同じ出典にもとづく。\nt{GLUT}バスと\nt{GABA}による流れの制御は確立しているが、ブロードキャストチャネルの役割とその協調は大部分が仮説である。

{\footnotesize
\setlength{\tabcolsep}{3pt}\renewcommand{\arraystretch}{1.15}
\begin{longtable}{>{\raggedright}p{0.5cm}>{\raggedright}p{4.0cm}>{\raggedright}p{5.6cm}>{\raggedright}p{2.3cm}>{\raggedright\arraybackslash}p{1.7cm}}
\toprule
No. & 主張 & 実験と測定されたもの & 出典 & 状態 \\
\midrule\endfirsthead
\toprule
No. & 主張 & 実験と測定されたもの & 出典 & 状態 \\
\midrule\endhead
1 & $8.6\cdot10^{10}$ 個のニューロンに対し、制御信号はわずか四種類~\eqref{eq:machine} & 脳のホモジネート中の細胞核の計数（等方性分画法）：成人男性でニューロンは $(86.1\pm8.1)\cdot10^9$ 個 & \cite{Azevedo2009} & 測定済み \\
2 & 命題~\ref{prop:architecture}の最初の二層：データは\nt{GLUT}のアドレスバスを流れ、結合の符号は送り手が決め、流れは\nt{GABA}ニューロンが方向づけ正規化する（第\ref{sec:neurontypes}, \ref{sec:gaba}章） & ニューロンあたり主要な伝達物質は一つ（測定の総説）；フィードフォワード抑制は錐体ニューロンが入力を加算する窓を狭める；全体の活動による除算は多くの領域で測定されている & \cite{Strata1999,Pouille2001,Carandini2012} & 測定済み \\
3 & 素子は信頼できない：シナプスはスパイクの大部分を失う（表~\ref{tab:computer}、第\ref{sec:code}章） & 海馬スライス、最小刺激：スパイクの半分以上がCA1で応答を起こさない。閾値や軸索伝導の失敗は除外され、原因は確率的な伝達物質放出 & \cite{AllenStevens1994} & 測定済み \\
4 & 脳は $\sim20$ ワットで動き、ニューロンあたり平均で毎秒約1スパイク（第\ref{sec:code}章）；このようなマシンをエンジニアはまだ作れていない & 測定されたコストからの推定~\eqref{eq:energy}：シナプスの仕事を含め1スパイクあたり約 $10^9$ 個のATP分子（げっ歯類の皮質）；皮質についての詳細な推定ではさらに低い発火率になる。技術の現状は総説による & \cite{Attwell2001,Lennie2003,MehonicKenyon2022} & 理論 \\
5 & 大多数のシナプスの学習は局所的な適格度トレースと一つの共通の数 $\delta$ の積（\eqref{eq:machine}の第2式、第\ref{sec:dofa}章） & サルの\nt{DOPA}ニューロンは報酬予測誤差に応答する；線条体スライスでは、ドーパミンがグルタミン酸入力の $0.3$--$2$\,s 後に来たときだけスパインを大きくする。トレースが測定された & \cite{Schultz1997,Yagishita2014} & 測定済み \\
6 & 出力ごとの専用の誤差配線があるのは運動学習の回路だけ（第\ref{sec:motorlearning}章） & 小脳：登上線維の信号と入力が同時に来るとその入力が弱まる；サルではプルキンエ細胞の複雑スパイクが運動誤差の方向を運び、補正を起こす & \cite{Ito2001,Herzfeld2018} & 測定済み \\
7 & 各ブロードキャストチャネルは数本の線の束（命題~\ref{prop:bundle}） & \nt{DOPA}について：同じ課題で異なるニューロンが報酬、運動、刺激の特徴を運ぶ；速い誤差と遅いドーパミンレベルは別物。\nt{SERO}, \nt{NORA}, \nt{ACET}ではこの分解はあまり調べられていない & \cite{Engelhard2019,Hamid2016,Mohebi2019} & 測定済み \\
8 & \nt{SERO}：レベル調整器であり、仮説では時間幅 $\tau_\gamma$ も決める（第\ref{sec:serocomputer}章） & 自己抑制：自らの5-HT$_{1A}$受容体のアゴニストが縫線核のセロトニンニューロンを抑制する（測定済み）。時間幅は理論で提案された；最も強い実験は、マウスでこれらのニューロンを光遺伝学的に活性化すると遅延報酬を待つ時間が延びること & \cite{Sprouse1987,Doya2002,Miyazaki2014} & 仮説 \\
9 & \nt{NORA}：ゲイン $g$ が探索か決定かを選ぶ（第\ref{sec:nora}章） & SN比の向上：覚醒したサルの聴覚皮質で、ノルアドレナリンは背景活動を同種の鳴き声への応答より強く抑える（測定済み）；乗算的ゲインはモデル；探索と決定の選択での役割は理論 & \cite{Foote1975NE,ServanSchreiber1990,AstonJones2005} & 仮説 \\
10 & \nt{ACET}：書き込みと読み出しを切り替える（第\ref{sec:acet}章） & 三つの作用（内部結合の弱化、入力の強化、可塑性の促進）は測定済み；モードの切り替えはヒトでのスコポラミン実験が間接的に支持 & \cite{HasselmoBower1992,Gil1997,Atri2004} & 仮説 \\
11 & 式~\eqref{eq:machine}はマシン全体を記述する & 第\ref{sec:glutcomputer}--\ref{sec:acet}章のモデルの要約で、各部分はそれぞれの章にもとづく；全体としては実験で検証されていない & 章内の導出 & 仮説 \\
12 & つまみは共通の制御なしに協調して働く：誤差、意外さ、書き込み、復帰（図~\ref{fig:timeline}） & 実験はない：模式図は定性的。個々の環は、\nt{NORA}と\nt{ACET}の分業の理論と、ヒトでの瞳孔と学習速度の関係 & \cite{YuDayan2005,Nassar2012} & 仮説 \\
13 & 一つの共通信号による学習は、結果に影響するニューロンが多いほど遅い（命題~\ref{prop:broadcastcost}） & 線形ネットワークの学習曲線の計算：出力の摂動による学習は勾配降下より出力の数の倍だけ遅い & \cite{Werfel2005} & 理論 \\
14 & 皮質が多層回路をどう調整するかは不明；候補はスパイクのバースト、ニューロンの枝の中での計算、予測による自己教師あり学習 & 誤差の担い手としてのバーストはモデル；皮質ニューロンの詳細モデルの応答を再現できたのは $5$--$8$ 層のネットワークだけ（モデル）；ヒト皮質ニューロンの枝で排他的論理和を生むカルシウムスパイクはスライスで測定済み；自己教師あり学習は証拠の総説 & \cite{Payeur2021,Beniaguev2021,Gidon2020,Keller2018} & 仮説 \\
\bottomrule
\end{longtable}}

\section{第\ref{sec:sensors}章　マシンの入力}

センサの仕組みは、単一細胞の電流記録、蝸牛の振動の記録、網膜の細胞の計数によって直接測定されている。感度の限界とショットノイズの式は物理から導かれる。センサの符号が情報を最大にするよう調整されているというのは、強い裏づけのある仮説である。

{\footnotesize
\setlength{\tabcolsep}{3pt}\renewcommand{\arraystretch}{1.15}
\begin{longtable}{>{\raggedright}p{0.5cm}>{\raggedright}p{4.0cm}>{\raggedright}p{5.6cm}>{\raggedright}p{2.3cm}>{\raggedright\arraybackslash}p{1.7cm}}
\toprule
No. & 主張 & 実験と測定されたもの & 出典 & 状態 \\
\midrule\endfirsthead
\toprule
No. & 主張 & 実験と測定されたもの & 出典 & 状態 \\
\midrule\endhead
1 & センサは変換器である：受容体が電流を生み、目と耳の感覚細胞の出力は\nt{GLUT}バスへのグルタミン酸である；最初の細胞はスパイクを出さない（図~\ref{fig:transducer}） & カメの桿体と錐体は興奮性アミノ酸を放出する：切り取った膜片のグルタミン酸チャネルがそれを捉える。ラット蝸牛の内有毛細胞：神経線維終末の電流はグルタミン酸AMPA受容体を通り、放出頻度は細胞のなめらかな脱分極とともに $150$\,Hz まで増す & \cite{CopenhagenJahr1989,GlowatzkiFuchs2002} & 測定済み \\
2 & 暗所の光センサは光子1個に約1ピコアンペアの電流で応答する & ヒキガエル桿体の外節への吸引電極：弱い閃光への応答は量子化されており、1事象は $\approx1$\,pA（暗電流の $5\%$）、ばらつき $0.2$\,pA、効率 $0.5$。ヒトは光子1個の到来を偶然より高い頻度で報告する & \cite{Baylor1979,Tinsley2016} & 測定済み \\
3 & 音のセンサは1ナノメートル未満の変位を検出する & メスバウアー法、モルモット蝸牛の基底膜の $16$--$19$\,kHz の位置：聴神経の応答しきい値で膜の速度は約 $0.04$\,mm/s、すなわち変位 $\approx0.35$\,nm（換算は本書による）。測ったのは膜であって、センサの毛束ではない & \cite{Sellick1982,Hudspeth1989} & 測定済み \\
4 & 暗所での精度は光子のショットノイズ~\eqref{eq:photonnoise}で制限される；弱い光では蓄積が長くなる & 式はポアソン統計から導かれる。桿体では弱い光のもとでのノイズが独立な量子事象の和と一致する；明るい背景では閃光への応答が小さく短くなる。「十分の数秒」という数値はここでは別に検証されていない & \cite{Baylor1979} & 理論 \\
5 & 入力のレンジは光で10桁、音で12桁だが、スパイクの発火率は3桁未満 & 音について：特徴周波数での基底膜の応答は $40$--$80$\,dB の帯域で傾き $0.2$\,dB/dB まで増し、圧縮は $100$\,dB を超えても見える。桁数そのものは標準的な推定で、章では出典なし & \cite{Ruggero1997} & 測定済み \\
6 & センサの応答は~\eqref{eq:nakarushton}で、$\sigma$ は平均の明るさを追い、曲線を対数軸に沿ってずらす（図~\ref{fig:adaptation}） & 式は最初、魚の網膜のS電位に当てはめられた。カメの錐体：応答は $V=I V_m/(I+\sigma)$ に従い、明るさの $\sim3.5$ 桁をカバーする；背景に $2$--$3$ 分さらすと、曲線は幅を保ったまま水平にずれる。全体の活動による除算との関係は総説 & \cite{NakaRushton1966,NormannPerlman1979,Carandini2012} & 測定済み \\
7 & センサは変化を伝え、その符号は1スパイクあたりの情報が最大になるよう調整されている（命題~\ref{prop:sensors}） & 原理は理論的に提案された。最も強い証拠：ハエの第1層ニューロンの「コントラスト $\to$ 応答」特性は、自然情景のコントラストの累積分布と一致する。こうして情報が最大になる & \cite{Barlow1961,Laughlin1981} & 仮説 \\
8 & オン型とオフ型のセンサは2線式符号である；イベントカメラはそれをシリコンで再現する & 実験の総説：網膜のオンチャネルとオフチャネルは最初のシナプスですでに分かれ、別々に遮断できる。カメラ：フレームなしの $128\times128$ 画素、レンジ $120$\,dB、遅延 $15$\,\textmu s & \cite{Schiller1992,Lichtsteiner2008,Gallego2022} & 測定済み \\
9 & 目には $\sim10^8$ 個の光センサと $\sim10^6$ 個の出力ニューロンがある：$\sim100$ 分の1への圧縮 & ヒト網膜の全載標本での計数：桿体は平均 $92\cdot10^6$ 個、錐体は $4.6\cdot10^6$ 個；出力（神経節）細胞は目によって $0.7\cdot10^6$ から $1.5\cdot10^6$ 個 & \cite{Curcio1990,CurcioAllen1990} & 測定済み \\
10 & マウスでは目の出力チャネルは30種類を超える & マウス：$11\,000$ 個を超える出力細胞の2光子カルシウムイメージングと応答のクラスタリング。機能的な型は $30$ を明らかに超える。ヒトでは数は測定されていない & \cite{Baden2016} & 測定済み \\
11 & モルモットの目は $\sim875\,000$ bit/s を送る；ヒトの目に換算すると $\sim10^7$ bit/s & モルモット、多電極アレイ、自然情景の中の運動：細胞あたり $6$--$13$ bit/s、$\sim10^5$ 個の出力細胞で $\sim875\,000$ bit/s。ヒト（$\sim10^6$ 個の細胞）の $\sim10^7$ bit/s は著者らの換算 & \cite{Koch2006} & 測定済み \\
12 & 耳はほぼ対数的な周波数マップを持つバンドパスフィルタのバンクである（図~\ref{fig:filterbank}） & 基底膜の振動が最大になる場所は周波数によって決まる；「周波数--場所」関数はほぼ指数関数的で、一つの形でヒトと生きた動物の蝸牛を記述する & \cite{Greenwood1990,Robles2001} & 測定済み \\
13 & 共振器は能動的である：一部のセンサが弱い振動を振幅で数千倍（$66$--$76$\,dB）増幅し、音が大きくなると増幅は下がる & チンチラ蝸牛基部のレーザー振動計測：特徴周波数の弱い純音で、アブミ骨に対する増幅は $66$--$76$\,dB；死後は感度が $60$--$81$\,dB（振幅で $10^3$--$10^4$ 倍）下がり、圧縮が消える。力の源は外有毛細胞 & \cite{Ruggero1997,Robles2001} & 測定済み \\
14 & 耳の増幅器は自励発振の境界に置かれている & 計算：ホップ分岐点の近くでは、レンジの圧縮、鋭い同調、結合音が避けられない。これは聴覚の既知の非線形性のすべてである。蝸牛がまさにそう調整されていることは直接には証明されていない & \cite{Eguiluz2000} & 仮説 \\
15 & 低い周波数ではスパイクは音と位相をそろえて出て（モルモットでは同期は $\sim600$\,Hz から弱まり、$\sim3.5$\,kHz まで認められる）、それから方向が求まる；高い周波数では場所符号が残る & モルモットの聴神経：位相同期は約 $600$\,Hz で落ち始め、$3.5$\,kHz を超えると消える；ネコとサルでは境界がほぼ1オクターブ高い。両耳への到達時間差は総説 & \cite{PalmerRussell1986,Grothe2010} & 測定済み \\
16 & その他のセンサ（表~\ref{tab:sensors}）：嗅覚は $\sim400$ 種類の受容体をニューロンあたり1種類持ち、組み合わせ符号を使う；味覚は一部ATPとセロトニンを使う；触覚は速く順応するセンサと遅く順応するセンサ；温と冷は別々 & マウス：一つの受容体が多くの匂いを認識し、一つの匂いが多くの受容体を活性化し、異なる匂いは異なる組を活性化する。ATP受容体がないと味覚神経は応答しない；味蕾はセロトニンを放出する。ヒトの手の皮膚の単一線維記録：順応の速さで四つの型。冷のチャネルは温のチャネルとは別物 & \cite{Mombaerts2004,Malnic1999,Finger2005,Huang2005,JohanssonVallbo1979,McKemy2002} & 測定済み \\
\bottomrule
\end{longtable}}

\section{第\ref{sec:recognition}章　認識}

認識の速さ、最初の段の仕組み、応答の線形な読み出し、時間の連続性による学習は、ヒトとサルで測定されている。連鎖全体を二つの演算に帰着させることと動画からの学習は本書のモデルで検証した。錯視の説明は、十分な根拠のあるものから論争中のものまである。

{\footnotesize
\setlength{\tabcolsep}{3pt}\renewcommand{\arraystretch}{1.15}
\begin{longtable}{>{\raggedright}p{0.5cm}>{\raggedright}p{4.0cm}>{\raggedright}p{5.6cm}>{\raggedright}p{2.3cm}>{\raggedright\arraybackslash}p{1.7cm}}
\toprule
No. & 主張 & 実験と測定されたもの & 出典 & 状態 \\
\midrule\endfirsthead
\toprule
No. & 主張 & 実験と測定されたもの & 出典 & 状態 \\
\midrule\endhead
1 & ずれがあると、画素ごとのテンプレート比較はコイン投げと変わらない；段の対~\eqref{eq:scstage}は図形をどの位置でも認識する & \texttt{models/\allowbreak recognition\_models.py}、$24$ 点の「網膜」、$4000$ 試行：反転した点の割合 $0$, $0.1$, $0.2$ でテンプレートは $0.49$, $0.51$, $0.52$；$S$ と $C$ の対は $1.00$, $0.86$, $0.70$ & 章内の導出 & モデル \\
2 & 絵の中の動物は $\sim150$\,ms で認識される。十数段を1段 $10$--$20$\,ms で1回通過する & ヒト、$20$\,ms 提示の写真で「動物がいるかいないか」の課題：二つの答えの誘発電位は $\sim150$\,ms で分かれる。サル：最初の応答の時刻は段を追って遅くなる（V1が最も早く、次にV2とV4）。段の数と「1段に0個か1個のスパイク」はこれらの数値からの推論 & \cite{Thorpe1996,Schmolesky1998} & 測定済み \\
3 & $S$ 段は部品の AND、$C$ 段は位置についての最大値（命題~\ref{prop:conveyor}） & ネコの一次視覚野：単純細胞は特定の場所の特定の向きのエッジに、複雑細胞は同じ向きにより広い領域で応答する。複雑細胞の細胞内記録：多くの細胞で2本のバーへの応答は二つの応答の大きいほうに近く、平均では max 演算に最も近い。相互抑制による最大値は命題~\ref{prop:wta}、全体の活動による除算は総説 & \cite{HubelWiesel1962,Lampl2004,Carandini2012} & 測定済み \\
4 & 連鎖を上るほど組み合わせは大きく複雑になり、応答は位置にますます依存しなくなる & サル、刺激を「臨界特徴」まで単純化：V4と後部下側頭皮質の細胞は小さな受容野で形、色、テクスチャの複雑な組み合わせに応答する；前部下側頭皮質では受容野が大きい。モデルの $S$ と $C$ の交互の繰り返しはこの描像を再現する & \cite{KobatakeTanaka1994,Riesenhuber1999} & 測定済み \\
5 & 畳み込みネットワークはこの交互の繰り返しを再現し、上位段の応答を最もよく予測する；物体はニューロン集団の発火率から線形に読み出せる & 脳へのあてはめなしに認識を学習したネットワーク：最上層はサル下側頭皮質のニューロン応答を、中間層はV4の応答を予測する。ランダムに選んだ $\sim100$ 個の細胞の $12.5$\,ms 以上の窓での活動から、線形分類器が位置や大きさが違っても物体とカテゴリーを認識する & \cite{Fukushima1980,Yamins2014,Hung2005} & 測定済み \\
6 & トレース付きヘッブ則~\eqref{eq:traceinvariance}は、トレースが $\sim20$\,ms 以上なら、ラベルなしの動画から不変性を学習する & 遅い特徴とトレース付き規則のアイデアは理論。\texttt{models/\allowbreak recognition\_models.py}、二つの $C$ ニューロン、入力 $16$、$10$ 回の実行、物体間の休止 $0.3$\,s。$\tau_{\text{tr}}=5$\,ms で図形との一致は平均 $0.52$、$10$\,ms で $0.56$；$20$, $50$, $200$\,ms ではすべての実行で $1.00$（$30$\,ms では10回中1回誤り） & \cite{Foldiak1991,WiskottSejnowski2002} & モデル \\
7 & 脳の中の不変性は時間の連続性から学習される & サル：視野の特定の場所への眼の跳躍中に物体をこっそり差し替えた；下側頭皮質ニューロンの位置不変性は差し替えに合わせて変わり、1時間後にはすでに有意だった。これが視覚学習の主要な規則だというのは仮説 & \cite{LiDiCarlo2008} & 測定済み \\
8 & 運動：方向検出器、次に広い領域についての同じ AND／OR の対 & ウサギ網膜の方向検出器；サルのMT野では記録した $110$ 個のニューロンすべてが方向選択的で、運動への応答はV1より強い & \cite{Barlow1965,Albright1984} & 測定済み \\
9 & 上位段は下へ予測を送り、上へは誤差が送られる & モデルは一次視覚野の受容野の古典外効果を説明する；個々の観察は整合する（総説）が、パイプラインの一般的な仕組みとしては証明されていない & \cite{RaoBallard1999,Keller2018} & 仮説 \\
10 & 難しい画像にはフィードバックと何周かの処理が要る & サル：「難しい」画像では下側頭皮質での決定が $\sim30$\,ms 遅れて形成される；この遅い応答は順方向ネットワークではうまく予測できず、フィードバックのあるネットワークやより深いネットワークのほうがよく予測する & \cite{Kar2019} & 測定済み \\
11 & 気づかれない画素の改ざんはネットワークをだます；ヒトの視覚ははるかに強いが無敵ではない & ネットワークで：特別に選んだ小さな摂動が答えを変える；「パンダ $\to$ テナガザル」の例は2番目の論文から。ヒトでは短時間提示のとき、ネットワーク間でよく転移する改ざんが選択をずらす & \cite{Szegedy2014,Goodfellow2015,Elsayed2018} & 測定済み \\
12 & 期待の錯視：信頼できない入力は期待に譲る。淡いものは「より遅く」動き、ドレスの見え方は人によって違う（\ref{sec:illusions}節） & コントラストの低い格子はより遅く動いて見える。$1401$ 人の調査：ドレスは三つの安定した見え方に分かれ、照明のヒントで色が切り替わる；$\sim13\,000$ 人の観察者では照明についての想定が決め手になる。遅い運動を期待するベイズモデルが一連の運動錯視を説明する & \cite{Thompson1982,LaferSousa2015,Wallisch2017,Weiss2002,Purves2011} & 仮説 \\
13 & ゲイン調整：滝、傾き、コントラスト；ヘルマン格子は近傍抑制ではない & ウサギ網膜の方向検出器は長い運動の間に発火率を下げ、停止後は背景以下に落ちる。傾きの錯視は近傍の活動による除算のモデルで再現される。網膜の出力ニューロンの応答を変えない格子の変形で斑点が消える。近傍抑制による説明は否定された & \cite{BarlowHill1963,Schwartz2009,SchillerCarvey2005} & 測定済み \\
14 & 遅延の補償：動く物体は閃光より前にあるように見える & 錯視は測定済み。心理物理は運動の前方延長とも遅延の差とも矛盾する；閃光後 $\sim80$\,ms の出来事にもとづく事後的な説明が提案されている。論争は決着していない & \cite{Nijhawan1994,EaglemanSejnowski2000} & 仮説 \\
15 & 静止した「ヘビ」が回る：遅延の差~\eqref{eq:latency}が方向検出器に読まれる & 錯視は色がなくても働く；隣り合う要素の対が単独で錯視を生む；方向選択的なサル皮質のニューロンは、この静止した対に方向性のある応答を示す。ヒトでは眼のマイクロサッカードと瞬きが回転を起動する & \cite{Conway2005,OteroMillan2012} & 測定済み \\
16 & 多義図形：相互抑制、疲労、ノイズ；切り替えは主にノイズが起こす & 競合するニューロン集団のモデル：切り替えを起こすのは\emph{ノイズ}であり、ノイズがなければ切り替えは起きない；刺激の強さとともに切り替え頻度が増すことを説明する。ここでは疲労の役割は小さい & \cite{MorenoBote2007} & 仮説 \\
17 & 錯視の一部はマシンが学習する課題から生じる & 動画の次フレームを予測するよう学習したネットワークは静止した輪の回転を「見て」、対照画像では見ない；画像のノイズ除去と鮮鋭化のネットワークはヒトと同様に色と明るさの錯視にかかる & \cite{Watanabe2018,GomezVilla2020} & 測定済み \\
\bottomrule
\end{longtable}}

\section{第\ref{sec:muscles}章　筋肉への出力}

出力の仕組み、すなわち運動単位、シナプスの安全率、大きさの順の動員、筋肉の対、伸張ループ、リズム発生器は直接測定されている。遅延ループの安定条件は定理である。体の内部モデルは十分な根拠のある仮説である。図の数値は本書のモデルによる計算である。

{\footnotesize
\setlength{\tabcolsep}{3pt}\renewcommand{\arraystretch}{1.15}
\begin{longtable}{>{\raggedright}p{0.5cm}>{\raggedright}p{4.0cm}>{\raggedright}p{5.6cm}>{\raggedright}p{2.3cm}>{\raggedright\arraybackslash}p{1.7cm}}
\toprule
No. & 主張 & 実験と測定されたもの & 出典 & 状態 \\
\midrule\endfirsthead
\toprule
No. & 主張 & 実験と測定されたもの & 出典 & 状態 \\
\midrule\endhead
1 & 運動単位：1個の\nt{ACET}ニューロンが数百から2000本ほどの線維のグループを制御する；細かな調整を要する筋肉では単位が小さい & ヒトの筋肉、運動軸索と筋線維の計数：ニューロンあたりの線維は平均で手の骨間筋で約 $340$、腕橈骨筋で約 $410$、腓腹筋内側頭で約 $1930$。最も小さな単位の正確な数は章では挙げていない。 & \cite{Feinstein1955mu} & 測定済み \\
2 & 筋肉上のシナプスは、線維の発火に必要な量の数倍の伝達物質を放出する & 神経筋シナプスでの測定の総説：安全率（放出された伝達物質と閾値の比）は成体哺乳類で通常 $3$--$5$；ヒトでは安全率は放出量よりも受け手の膜のひだによって保たれている。 & \cite{WoodSlater2001} & 測定済み \\
3 & 単収縮~\eqref{eq:twitch}は $50\ms$ 程度の時間 $T_c$ で立ち上がる & 随意収縮中のヒトの手の単一運動単位、単位のスパイクで加算平均した力：立ち上がり時間 $30$--$100\ms$、$80\,\%$ の単位で $70\ms$ 未満。関数 $(t/T_c)e^{1-t/T_c}$ は運動単位プールのモデルからの近似。 & \cite{MilnerBrown1973rec, Fuglevand1993} & 測定済み \\
4 & 単収縮の和~\eqref{eq:forcesum}：$5$, $15$, $40$ スパイク/s で力は $0.2$--$1.1F_1$、$1.8$--$2.2F_1$、約 $5.4F_1$（図~\ref{fig:tetanus}） & \texttt{models/\allowbreak muscle\_\allowbreak models.py} による計算、$T_c=50\ms$：最後の $0.3\,\text{s}$ で $0.21$--$1.10$、$1.76$--$2.19$、$5.28$--$5.49\,F_1$。線形の加算は単純化：実際の筋肉の飽和はモデルにない。 & \texttt{models/\allowbreak muscle\_\allowbreak models.py} & モデル \\
5 & 単位は大きさの順に動員される；最も弱いものは最も強いものの100分の1 & ネコの前根：反射入力を強めていくと、前根でのスパイクが小さいニューロン、つまり小さいニューロンが最初に発火する。ヒトの手の骨間筋：単位の単収縮の力は $0.1$ から $10$~g で、単位が動員される力にほぼ比例して増える。 & \cite{Henneman1957, MilnerBrown1973rec} & 測定済み \\
6 & 力の刻みは力に比例する、$\Delta F/F\approx q-1$~\eqref{eq:sizeprinciple}：浮動小数点 & 章で力の等比数列から導出。実験と整合する：大きな力では、同じ力の増分に対して新たに動員される単位がずっと少ない。 & 章内の導出；\cite{MilnerBrown1973rec} & 理論 \\
7 & 力のばらつきは力そのものに比例して増える & 母指の筋の随意等尺性収縮：力の標準偏差は平均の力に比例して増える；電気刺激で同じ力を出させるとばらつきは一定。プールのモデル：力の順の動員が線形の増加を生む。 & \cite{Jones2002sdn} & 測定済み \\
8 & 動きのなめらかさは信号依存ノイズの帰結である & モデル：このようなノイズのもとで終点のばらつきを最小にする軌道は、測定された眼と腕の軌道を再現する。この原因を他の原因と区別する直接の実験はない。 & \cite{HarrisWolpert1998} & 仮説 \\
9 & 筋肉の対の力の差がトルクを、和が剛性を決める~\eqref{eq:antagonists} & 同時収縮によるインピーダンス制御の理論。ヒトの腕：小さな摂動で測った各関節の剛性は、その関節のトルクにほぼ比例する；同時収縮の比率を変えると手先の剛性楕円の形と向きが変わる。 & \cite{Hogan1984, GomiOsu1998stiff} & 測定済み \\
10 & 伸張センサは自分の\nt{ACET}ニューロンを興奮させ、抑制性の介在ニューロンを通じて拮抗筋をオフにする（図~\ref{fig:antagonists}） & ネコの脊髄：伸張センサの線維に興奮させられる単一の介在ニューロンが、拮抗筋の\nt{ACET}ニューロンに単シナプス性の抑制性電位を生む。シナプス遅延 $0.28$--$0.42\ms$。介在ニューロンの伝達物質（グリシン）はこの実験では同定されていない。 & \cite{Jankowska1972Ia} & 測定済み \\
11 & ループの遅延：脊髄 $25$--$30\ms$、皮質経由はその1.5〜3倍、目経由 $100$--$200\ms$ & ヒトの腕、関節への摂動：筋肉の短潜時応答は $20$--$45\ms$、長潜時応答は $45$--$105\ms$、随意応答は $120$--$180\ms$。運動中に見える指の位置をずらすと、修正は平均 $160\ms$ 後。 & \cite{Pruszynski2008rapid, SaundersKnill2003vis} & 測定済み \\
12 & $\dd x/\dd t=-Gx(t-d)$ は $Gd<\pi/2$~\eqref{eq:delaystable}で安定；境界での周波数は $1/(4d)$ & 遅延方程式の根についての定理。\texttt{models/\allowbreak muscle\_\allowbreak models.py} による検証、$d=30\ms$：$3\,\text{s}$ の時点で振幅は $G=40$ で $3\cdot10^{-6}$、$G=50$ で $0.13$、$G=55$ で $38$（単位は毎秒、境界は $52$）。 & \cite{Hayes1950} & 理論 \\
13 & 生理的振戦 $8$--$12$~Hz；伸張ループはその源の一つ & 測定の総説：健康な人には約 $10$~Hz 帯の細かな震えがある；その起源は混合的で、中枢の振動、単位の発火の性質、機械的共振、反射ループの共振があり、条件によって優勢なものが異なる。 & \cite{McAuleyMarsden2000} & 仮説 \\
14 & 脳は体の内部モデルによって指令の結果を予測する & ヒトが物体をつまんで持ち、腕を動かす：慣性、粘性、弾性の負荷のいずれでも、把持力はループの遅れなしに負荷と同時に変わる。つまり負荷は予測されている。総説がこうした実験をまとめる；ニューロン内のモデルそのものの直接観察はない。 & \cite{FlanaganWing1997grip, WolpertGhahramani2000} & 仮説 \\
15 & 歩行、呼吸、咀嚼のリズムは、一定の入力のもとで局所的なネットワークが生む & 実験の総説：摘出した神経系（脊髄、神経節）は、リズミカルな入力も筋肉からのフィードバックもなしにリズミカルな運動パターンを出す。 & \cite{MarderBucher2001} & 測定済み \\
16 & 疲労をともなう相互抑制~\eqref{eq:halfcentre}は周期 $0.84\,\text{s}\approx2\tau_a$ のリズムを生む（図~\ref{fig:halfcentre}） & 定理：相互抑制と適応を持ち、一定の入力のもとで安定状態を持たないネットワークは必ず振動する。\texttt{models/\allowbreak muscle\_\allowbreak models.py} による計算（$I=1$, $\beta=2$, $b=2$, $\tau=20\ms$, $\tau_a=0.4\,\text{s}$）：周期 $0.84\,\text{s}$。実際の発生器がまさにこのような構成であることは示されていない。 & \cite{Matsuoka1985osc} & モデル \\
\bottomrule
\end{longtable}}

\section{第\ref{sec:pain}章　痛み}

痛みのセンサ、2本の配線、脊髄のゲート、下行性の音量つまみ、感作、注意の奪取は直接測定されている。ゲートと感作の式は本書の単純化したモデルである。マシンが警報の音量を選ぶ規則は仮説である。

{\footnotesize
\setlength{\tabcolsep}{3pt}\renewcommand{\arraystretch}{1.15}
\begin{longtable}{>{\raggedright}p{0.5cm}>{\raggedright}p{4.0cm}>{\raggedright}p{5.6cm}>{\raggedright}p{2.3cm}>{\raggedright\arraybackslash}p{1.7cm}}
\toprule
No. & 主張 & 実験と測定されたもの & 出典 & 状態 \\
\midrule\endfirsthead
\toprule
No. & 主張 & 実験と測定されたもの & 出典 & 状態 \\
\midrule\endhead
1 & 高い閾値：痛みのセンサの加熱閾値は $41^\circ$ から $46$--$48^\circ$ & 単一線維記録。サルの皮膚：遅い（C）センサの閾値の中央値は $41^\circ$、速いセンサの第2型は $46^\circ$、第1型は $53^\circ$ 超。ヒトの皮膚：圧力にも応答するCセンサは $40.7^\circ$、圧力に応答しないものは $48^\circ$。 & \cite{Treede1995heat, Weidner1999cfib} & 測定済み \\
2 & 痛みのセンサは触覚のセンサよりはるかに順応しにくく、軽い損傷の後にはより敏感になる；強い加熱の後の低下は一つの型で測定 & 皮膚の軽いやけど（$50^\circ$、$100\,\text{s}$）：$5$--$10$~分後、ヒトの痛みの閾値とサルのCセンサの閾値が $2$--$6^\circ$ 下がる；他の皮膚センサにはこのずれはない。速いセンサの第2型では、強い加熱の後に応答が抑えられる。触覚センサとの順応の比較は一般的な推定で、ここでは個別の実験で裏づけられていない。 & \cite{LaMotte1982hyper, Treede1995heat} & 測定済み \\
3 & 2本の配線：速いほうは $5$--$30$~m/s、遅いほうは $0.5$--$2$~m/s；到着時間 $t=L/v$~\eqref{eq:paintime} & 伝導速度：サルの速い痛みのセンサは平均 $15$ と $25$~m/s（二つの型）；ヒトのCセンサは $0.8$--$1.0$~m/s。$L=1$~m ならそれぞれ数十ミリ秒と約1秒。 & \cite{Treede1995heat, Weidner1999cfib} & 測定済み \\
4 & 痛みは2回来る：最初は速く正確に、次に鈍く長く & ヒト前腕の加熱に対する反応時間：温度の上昇とともに $1100$ から $400\ms$；有毛部の皮膚の強い加熱は二重の痛みを生むが、手のひらでは生まない。第一の痛みを運べるのは $6$~m/s より速い線維だけで、潜時からそれに対応するのは手のひらにない速いセンサの第2型である。役割の分担（「どこ」と「どれほど悪いか」）は解釈。 & \cite{CampbellLaMotte1983first, Treede1995heat} & 測定済み \\
5 & ゲート：脊髄の出力ニューロンは痛みと触覚を受け、抑制性の介在ニューロンは触覚で興奮する（図~\ref{fig:gate}） & ゲートの構成は理論として提案された。マウス、遺伝的標識とニューロン群の除去：ソマトスタチンを持つ興奮性ニューロンは機械的な痛みに必要；ダイノルフィンを持つ抑制性ニューロンは触覚センサからの入力がそこへ届くのを防ぐ。これらのニューロンがないと軽い接触が痛みを起こす。 & \cite{MelzackWall1965, Duan2014} & 測定済み \\
6 & 触覚は痛みを弱める & ヒト、腕へのレーザー痛み刺激：同時の触覚は痛みの評定とパルスのエネルギーの識別能力を下げる；効果は、同じ皮膚分節の中で触れる点と痛む点の距離とともに弱まる。 & \cite{Mancini2014touch} & 測定済み \\
7 & ゲート理論の仮定：強い痛みは自ら抑制性介在ニューロンをオフにし、ゲートが開け放たれる & 章では元の理論の仮定と明記。マウスでの研究は、ゲートが触覚を痛みのチャネルに入れない仕組みを記述するが、痛みによる介在ニューロンのオフは示していない。 & \cite{MelzackWall1965} & 仮説 \\
8 & ゲート~\eqref{eq:gate}：閾値は触覚なしで $0.18$、触覚ありで $0.86$、$g=2$ で $0.11$；$\nu=1$ で触覚は出力を $1$ から $0.25$ に下げ、$\nu=2$ では効かない；触覚単独では通らない（図~\ref{fig:gatecurves}） & 本文の重みでの \texttt{models/\allowbreak pain\_\allowbreak gate.py} による計算が、これらの数値をすべて再現する。 & \texttt{models/\allowbreak pain\_\allowbreak gate.py} & モデル \\
9 & 脳幹からの下行路は、ゲートのゲインをどちらの向きにも変える & ラット延髄、尾の加熱中の記録：あるニューロンは引っ込めの前に発火し（「オン」）、別のニューロンは沈黙する（「オフ」）；この領域の弱い刺激は反射を抑える。総説：同じ回路が痛みの伝達を促進も抑制もし、オピオイドはそれを介して働く。 & \cite{Fields1983rvm, Fields2004} & 測定済み \\
10 & 楽になるという期待は、オピオイドのチャネルを通じて痛みを弱める & 急性の痛みを持つ人：楽になるという期待で痛みが減った人では、オピオイド受容体の遮断で痛みが他の人と同じレベルに戻った；それ以外の人では遮断は痛みを変えなかった。 & \cite{Levine1978} & 測定済み \\
11 & 脳は割り込みの損得によって警報の音量を選ぶ & 音量を選ぶ規則を測定した実験はない。 & --- & 仮説 \\
12 & 感作：損傷の後、ゲートの出力は増し、閾値は下がり、その一部は脊髄による；損傷刺激が止んだ後も続く & ラット：皮膚の損傷の後、屈曲反射の閾値が下がり、応答が増す；電気生理学的解析によれば、増大の一部は脊髄そのもので生じる。損傷刺激は刺激そのものより長く続く感度の変化を起こす。正確な回復時間は章では挙げていない。 & \cite{Woolf1983} & 測定済み \\
13 & 感作~\eqref{eq:sensitization}は正帰還である；ループが強いと、警報は治癒後もラッチしうる & モデル~\eqref{eq:sensitization}から導かれる（章末の問題）。治癒後も続く痛みがまさにこの形のラッチしたループであることを示す実験はない。 & 章内の導出 & 仮説 \\
14 & 痛みは負の報酬であり、それによる学習は時間差分誤差に従う & fMRI、ヒト、痛みの前触れの連鎖：腹側線条体と前部島皮質の活動が時間差分学習の信号と一致する。サル：\nt{DOPA}ニューロンの一部は不快なエアパフとその前触れで興奮し、別の一部は抑制される。 & \cite{Seymour2004, Matsumoto2009} & 測定済み \\
15 & 痛みはいまの作業に割り込む & ヒト、電気的な痛み刺激と音のプローブ：痛みの開始 $100\ms$ 後のプローブへの応答は明らかに遅れ、$1.5\,\text{s}$ 後には対照との差はもうない。総説はこうした実験を注意の奪取のモデルにまとめる。 & \cite{Crombez1996disrupt, EcclestonCrombez1999} & 測定済み \\
16 & 引っ込め反射は脊髄で閉じる & ラット：後角深層のニューロンは個々の筋肉の引っ込め反射の空間パターンを再現する；頸部からは逆行性に興奮させられず、脊髄の介在ニューロンである。 & \cite{Schouenborg1995wr} & 測定済み \\
\bottomrule
\end{longtable}}

\section{第\ref{sec:motorlearning}章　運動学習}

最小二乗則と専用の誤差配線の利点は定理である。誤差の配線を持つ回路の構成、一致による弱化、遅延へのトレースの調整、後効果と技能の自発的回復は測定されている。回路全体が教師あり学習として働き、内部モデルを作るというのは有力な仮説である。二過程モデルの数値は本書のモデルによる計算である。

{\footnotesize
\setlength{\tabcolsep}{3pt}\renewcommand{\arraystretch}{1.15}
\begin{longtable}{>{\raggedright}p{0.5cm}>{\raggedright}p{4.0cm}>{\raggedright}p{5.6cm}>{\raggedright}p{2.3cm}>{\raggedright\arraybackslash}p{1.7cm}}
\toprule
No. & 主張 & 実験と測定されたもの & 出典 & 状態 \\
\midrule\endfirsthead
\toprule
No. & 主張 & 実験と測定されたもの & 出典 & 状態 \\
\midrule\endhead
1 & 規則~\eqref{eq:lms}は平均すると二乗誤差の勾配をたどる & 適応フィルタ理論の結果：入力と出力の誤差に比例する補正は、平均二乗誤差の確率的勾配降下である。 & \cite{WidrowHoff1960} & 理論 \\
2 & 出力ごとに誤差の配線があれば速さは出力の数によらない；一つの共通の数では、ノイズを出すニューロンの数とともに落ちる（命題~\ref{prop:teachers}） & 線形ネットワークの学習曲線：ニューロンのランダムな摂動による学習は、摂動を受けるニューロンの数の倍だけ正確な勾配より遅い。 & \cite{Werfel2005} & 理論 \\
3 & 誤差の配線を持つ回路は教師あり学習として働く（命題~\ref{prop:errorwire}） & 回路の構造から導かれた理論。7--10行の実験はそれと整合する；回路全体がまさにそう働くことは証明されていない。 & \cite{Marr1969, Albus1971} & 仮説 \\
4 & 拡張層：約 $7\cdot10^{10}$ 個の小さな\nt{GLUT}ニューロン、脳の残り全部より多い & 溶解した組織の懸濁液での細胞核の計数：ヒトの小脳には脳全体の $86\cdot10^9$ のうち $69\cdot10^9$ のニューロン。立体解析学：高齢男性の小脳に $101\cdot10^9$ 個の小さなニューロン。 & \cite{Azevedo2009, Andersen1992cereb} & 測定済み \\
5 & 入力の次元を上げると、望みの関係が線形になる & 分割の数についての定理：$n$ 個の入力の重み付き和としきい値は、一般の位置にある約 $2n$ 個のベクトルの任意のラベル付けを実現する。小脳がまさにこれを使っているというのは3行目の仮説の一部。 & \cite{Cover1965} & 理論 \\
6 & 約 $3\cdot10^7$ 個の出力\nt{GABA}ニューロン、それぞれ $10^5$ 程度の入力 & ヒト小脳の立体解析学：出力ニューロンは $30.5\cdot10^6$ 個。電子顕微鏡、ラット：出力ニューロン1個あたり拡張層からの入力は約 $1.75\cdot10^5$。出力ニューロンの抑制性伝達物質は総説による。 & \cite{Andersen1992cereb, NapperHarvey1988, Ito2001} & 測定済み \\
7 & 出力ニューロンあたり誤差の配線はちょうど1本、毎秒1回ほど、毎回強力なバースト & 記録の総説：各出力ニューロンは1本の登上\nt{GLUT}入力を受け、それが $1$~Hz 程度の頻度で複雑スパイクを起こす。 & \cite{Ito2001} & 測定済み \\
8 & バーストの確率は運動誤差の方向に依存する & サル、小脳の眼球運動領域、サッカード適応：視覚誤差の方向が複雑スパイクの確率を、大きさがそのタイミングを決める；バーストは次の試行の単純スパイクを変え、細胞の好みの誤差の方向に運動をずらす。 & \cite{Herzfeld2018} & 測定済み \\
9 & 誤差のバーストと活動中の入力が一致すると、その入力が弱まる（$-\eta e_i x_j$） & 実験の総説：拡張層からの入力と誤差の配線を同時に刺激すると、その入力が長期的に弱まる；入力だけの刺激では弱まらない。 & \cite{Ito2001} & 測定済み \\
10 & トレースの長さは誤差の遅延に合わせてある：遅延 $100\ms$ では $100\ms$ 前に活動していた入力が最も強く弱まる & スライスと生きたマウス：眼球運動学習を担う領域では、弱化は入力と誤差の間隔約 $120\ms$（この課題で誤差信号が届くのにかかる時間）に鋭く同調している；別の領域では細胞ごとに異なる間隔に同調している。 & \cite{Suvrathan2016} & 測定済み \\
11 & 回路は体の内部モデルを学習する適応フィルタ~\eqref{eq:adaptivefilter}である & 回路の構造から導かれたモデル。学習したフィルタを体の伝達関数として測定した直接の実験はない。 & \cite{Fujita1982} & 仮説 \\
12 & 予測は自分の動きの結果を差し引くので、自分をくすぐることはできない & ヒト：自分で触れると、他人の同じ接触よりくすぐったくない；fMRIでは体性感覚皮質の自分の接触への応答が弱く、動きが皮膚に触れるとき小脳の活動が少ない。予測の差し引きによる説明は仮説。 & \cite{Blakemore1998} & 仮説 \\
13 & 外力のもとで外れは小さくなり、外力を取り除くと手は反対側に外れる & ヒトが力場の中でロボットのハンドルを持って手を伸ばす：軌道はまっすぐに戻る；力場を取り除いた後の後効果は最初のゆがみのほぼ鏡像で、訓練していない空間の部分にも転移する。 & \cite{ShadmehrMussaIvaldi1994} & 測定済み \\
14 & 二つの過程が学習する~\eqref{eq:tworate}；逆向きの外乱の後、技能はひとりでに戻る & ヒト、力場、次に短い逆向きの力場と誤差を固定した試行：補正はひとりでに古い値へ戻る；二過程モデルはこれを再現し、一過程モデルは再現しない。 & \cite{Smith2006} & 測定済み \\
15 & 図~\ref{fig:tworate}の数値：$200$ 回の動きで $0.84$（遅い過程 $0.63$）；逆向き $6$ 回；速い $-0.53$、遅い $0.46$；$40$ 回で $0.37$ まで回復 & \texttt{models/\allowbreak motor\_\allowbreak learning.py} による計算、$A_f=0.92$, $B_f=0.1$, $A_s=0.996$, $B_s=0.02$：$0.840$（$0.634$ と $0.205$）、$6$ 回、$-0.531$ と $0.457$、$40$ 回目でピーク $0.371$。 & \texttt{models/\allowbreak motor\_\allowbreak learning.py} & モデル \\
16 & 体が異なる速さで変わるので、二つの過程が必要になる & 理論：寿命の異なる複数の外乱のもとでの最適推定は、時定数の異なる複数のフィルタを与え、自発的回復と再学習の加速を説明する。 & \cite{Kording2007} & 仮説 \\
\bottomrule
\end{longtable}}

\section{第\ref{sec:chemoutput}章　化学出力}

心臓への2本の配線とその速さの違い、ホルモンカスケードの負帰還、分量の頻度による符号、概日発振器とその光による同期は直接測定されている。境界$K<8$は本章で導いた制御理論の定理である。カスケードのフィードバックそのものが脈動を生むというのは仮説であり、それを強く支持する実験がある。

{\footnotesize
\setlength{\tabcolsep}{3pt}\renewcommand{\arraystretch}{1.15}
\begin{longtable}{>{\raggedright}p{0.5cm}>{\raggedright}p{4.0cm}>{\raggedright}p{5.6cm}>{\raggedright}p{2.3cm}>{\raggedright\arraybackslash}p{1.7cm}}
\toprule
№ & 主張 & 実験と測定内容 & 出典 & 状態 \\
\midrule\endfirsthead
\toprule
№ & 主張 & 実験と測定内容 & 出典 & 状態 \\
\midrule\endhead
1 & 心臓のペースメーカーは単独で毎分約$100$回拍動する。安静時はブレーキが踏まれており、心拍数はふつう$60$--$70$程度 & ヒト、心臓への2本の配線をともに完全に遮断：固有心拍数は安静時心拍数より高く、加齢とともに下がり、成人で毎分約$100$拍。したがって安静時はブレーキが優位である。安静時心拍数$60$--$70$は典型値で、個別には引用していない。 & \cite{JoseCollison1970ihr} & 測定済み \\
2 & アクセル\nt{NORA}とブレーキ\nt{ACET}はローパスフィルタで、ブレーキの方が速い~\eqref{eq:gasbrake} & イヌ、心臓の迷走神経または交感神経を周波数変調刺激し、心房レートへの伝達関数を測定：どちらの経路もローパスフィルタで、交感神経側は遮断周波数がはるかに低く、約$1.7\,\text{s}$の純粋な遅れが加わる。特性は平均緊張度に依存するので、線形式~\eqref{eq:gasbrake}は近似である。 & \cite{Berger1989} & 測定済み \\
3 & ブレーキが速いのは、\nt{ACET}がセカンドメッセンジャーなしでカリウムチャネルを開き、\nt{NORA}は反応の連鎖を通じて作用するから & 心房細胞の膜パッチ：アセチルコリンで開くカリウムチャネルは、受容体と共役したGタンパク質の$\beta\gamma$サブユニットが、細胞内のセカンドメッセンジャーを介さずに開く。cAMPを介する\nt{NORA}の経路はここでは引用していない。 & \cite{Logothetis1987gbg} & 測定済み \\
4 & 血液はおよそ1分で全身を一周する。血中の\nt{ADRE}は受容体をもつすべての臓器に届く & 一般的な桁の見積もり。本章でもそのように述べており、出典は引用していない。 & --- & 推定 \\
5 & ホルモンカスケードは負帰還で囲まれている：最終段のホルモンが前段を抑える & 実験の総説：副腎皮質ホルモンは、下垂体からのACTH放出と視床下部からの放出ホルモン放出を、速い、中間、遅いの三つの時間域で抑制する。 & \cite{KellerWood1984fb} & 測定済み \\
6 & 同一の3段は$K<8$で安定~\eqref{eq:cascadestable}、境界での周期は$2\pi\tau/\sqrt3\approx3.6\tau$ & 本章で導出：周波数$\sqrt3/\tau$で3段は位相を$180^\circ$ずらし、$8$分の1に減衰させる。 & 本章で導出 & 理論 \\
7 & レベルの誤差は$1/(1+K)$で、この調節器はおよそ$10\,\%$より高い精度を保てない（命題~\ref{prop:cascade}） & 比例フィードバックに対する6行目の帰結。実際の軸は複数の段と複数の時間域にフィードバックをもつ（5行目）ので、この境界はモデル~\eqref{eq:cascade}についてのものであり、測定値ではない。 & 本章で導出 & 理論 \\
8 & $K=4$と$7$ではレベルが落ち着き、$K=12$では周期$3.7$--$3.9\,\tau$の規則正しい脈動（図~\ref{fig:cascade}） & \texttt{models/\allowbreak chem\_\allowbreak models.py}による計算：$K=12$で連続する周期は$3.9$、$3.7$、$3.8$、$3.7\,\tau$。$K=4$では計算終了時の振幅$0.003$。 & \texttt{models/\allowbreak chem\_\allowbreak models.py} & モデル \\
9 & ストレスホルモンは約1時間に1回のパルスで放出される。パルスはフィードバックそのものが生み、独立した発振器はない & ラット、放出ホルモンの持続注入：ACTHとコルチコステロンは生理的なほぼ1時間の頻度で振動する。パルスに視床下部からのリズム入力は必要ない。遅れのあるフィードバックをもつ下垂体--副腎モデルがこの振動を再現する。 & \cite{Walker2012glc, Walker2010} & 仮説 \\
10 & 受け手はレベルではなく分量の頻度を読む & 性腺刺激ホルモン放出ホルモンを自ら分泌できないサル：持続注入では下垂体ホルモンの放出は回復せず、1時間に1回の分量投与で回復する。持続注入に戻すと応答は再び消える。ハイパスフィルタとしての受容体の疲労は解釈である。 & \cite{Belchetz1978} & 測定済み \\
11 & 分量の頻度が違えば、同じ腺の異なる細胞への作用も違う & 同じサル：分量を1時間に$2$--$5$回に増やすと下垂体ホルモンが両方とも下がる。$3$時間に1回に減らすと一方（LH）が下がり他方（FSH）が上がるので、両者の比は頻度で決まる。 & \cite{Wildt1981gnrh} & 測定済み \\
12 & 概日リズムは数時間の遅れをもつ細胞内の負帰還が生む & 総説：時計遺伝子のタンパク質が遅れて自らの遺伝子の転写を抑える。転写と翻訳のループが約1日の周期を生む。 & \cite{TakahashiJS2017} & 測定済み \\
13 & 主発振器は数万個のニューロンからなる群で、体の他の部分を同期させる & 総説：視交叉上核が主概日発振器である。培養した個々のニューロンは独立した発振器で、ネットワークがそれらを同期させる。マウスでは片側に約$10^4$個のニューロン。 & \cite{Welsh2010scn} & 測定済み \\
14 & ヒトの固有周期は$24.18$時間 & 薄暗い光のもと、1日から切り離したスケジュールで暮らすヒト：メラトニン、体温、コルチゾールのリズムの周期は若年者と高齢者で平均$24.18$時間、ばらつきは小さい。 & \cite{Czeisler1999} & 測定済み \\
15 & 光は位相同期ループのように発振器を合わせる~\eqref{eq:pll}。必要な補正は1日約$11$~min & ヒト、1日のさまざまな時刻に$6.7$時間の明るい光を1回照射：振幅$5$時間のタイプ1位相反応曲線。体温最低点より前では遅れ、後では進む。式~\eqref{eq:pll}は標準モデル、$11$~min $=0.18$時間は算術である。 & \cite{Khalsa2003prc} & 測定済み \\
16 & 光がなければ同期はなく、リズムは固有周期で進む & 光をまったく知覚できない全盲者で、通常の社会生活を送る人：約半数でメラトニンとコルチゾールのリズムが1日と一致しない固有周期で進む。 & \cite{Sack1992blind} & 測定済み \\
\bottomrule
\end{longtable}}

\section{第~\ref{sec:debugging}~章　マシンのデバッグ}

電極に何が聞こえるか、活動の低次元性、硬い配列によるインターフェースの性能、脳とデコーダの共同学習、刺激による視覚上の点は、査読付きの実験で測定されている。データ量の見積もりは本章の算術である。ヒトに埋め込んだNeuralinkの装置については、確認できた限り、データを公表しているのは主に同社自身である。

{\footnotesize
\setlength{\tabcolsep}{3pt}\renewcommand{\arraystretch}{1.15}
\begin{longtable}{>{\raggedright}p{0.5cm}>{\raggedright}p{4.0cm}>{\raggedright}p{5.6cm}>{\raggedright}p{2.3cm}>{\raggedright\arraybackslash}p{1.7cm}}
\toprule
№ & 主張 & 実験と測定内容 & 出典 & 状態 \\
\midrule\endfirsthead
\toprule
№ & 主張 & 実験と測定内容 & 出典 & 状態 \\
\midrule\endhead
1 & 電極には半径約100マイクロメートル以内のニューロン、ふつう1個から数個のスパイクが聞こえ、波形で区別できる & ラットCA1野、細胞内と細胞外の同時記録：細胞外スパイクの波形は細胞内スパイクの幅と振幅を反映する。見積もりではテトロードは$60$--$100$個のニューロンを区別できるはずだが、実際に分離できるのは約6個。 & \cite{Henze2000extra} & 測定済み \\
2 & 脳には$8.6\cdot10^{10}$個のニューロンがあり、プローブが聞くのはおよそ1億分の1 & 溶解した組織の懸濁液中の核の計数：ヒトの脳のニューロンは$86\cdot10^9$個。割合は本章の算術。 & \cite{Azevedo2009} & 測定済み \\
3 & 運動皮質の数百個のニューロンの活動は10--20個の共通変数に収まる & サル運動皮質の記録の総説：集団活動は、多くのニューロンに共通する少数の潜在変数で記述できる。 & \cite{Gallego2017} & 測定済み \\
4 & 隣接ニューロンのノイズは相関しており、100個のニューロンは1000個とほぼ同じだけ運ぶ（命題~\ref{prop:pooling}） & サル視覚皮質：隣接ニューロンのノイズ相関係数は約$0.1$で、平均化の利得は100個のニューロンで飽和する。情報制限相関の理論。 & \cite{Zohary1994, MorenoBote2014} & 測定済み \\
5 & 生データ$2\cdot10^8$~bit/sに対しスパイクでは$8\cdot10^4$~bit/s~\eqref{eq:probedata} & スパイク列の容量~\eqref{eq:spikeentropy}による本章の算術。デジタル化のパラメータは実際のもの：Neuralinkのチップは$19.3$~kHz、チャネルあたり$10$~bit。 & 本章で導出；\cite{Rieke1997, Musk2019} & 理論 \\
6 & Neuralinkの最初のシステム：チップは増幅とデジタル化を行い、生データはケーブルで送られ、スパイクは外部のプログラムが検出する。チップ上の検出、密閉筐体、無線、ワイヤレス給電は、同社によればヒト用の装置のもの & 同社の論文：生データ全体をUSB-Cケーブルで送り、チップ外のプログラムがリアルタイムでスパイクを検出する。搭載した圧縮、ワイヤレス給電と伝送は臨床用装置の計画とされている。ヒトに埋め込んだ装置については同社の発表のみ。チップ上でのスパイク検出が必要だというのは、\eqref{eq:probedata}からの本章の結論。 & \cite{Musk2019}；Neuralinkの報告、査読付き論文なし & 測定済み（同社の論文）。ヒトについては同社の報告 \\
7 & $96$本のスレッドに$32$個ずつ最大$3072$個の電極。スレッドはロボットが血管を避けて挿入する & 同社の論文：$96$本のスレッドに$32$個ずつ$3072$個の電極。ロボットは毎分$6$本（$192$電極）を挿入する。配列を長期埋め込みしたラットで最大$70\,\%$の電極がスパイクを聞き取る。 & \cite{Musk2019} & 測定済み（同社の論文） \\
8 & ヒトでは$64$本のスレッドで約1000チャネル。一部のスレッドがずれた。カーソルは約$10$~bit/s & 同社の発表のみ。PubMedの検索ではヒトでの結果を示す査読付き論文は見つからなかった。本文中の但し書きと一致する。 & Neuralinkの報告。査読付き論文なし & 測定済み（同社の報告） \\
9 & 100電極の配列によるインターフェースがカーソルを操作する & 麻痺のある人、一次運動皮質に$96$電極：受傷3年後でも腕を動かす意図がスパイクを変える。デコーダは「ニューラルカーソル」（メール、テレビ）、義手と ロボットアームの操作を可能にする。 & \cite{Hochberg2006} & 測定済み \\
10 & 「想像上の手」による書字：毎分約$90$文字、$8$~bit/s未満 & 手が麻痺した人が手書きを試み、再帰型デコーダを使用：リアルタイムで毎分$90$文字、精度$94.1\,\%$、自動修正後は$99\,\%$超。ビットでの見積もりは本章の算術。 & \cite{Willett2021} & 測定済み \\
11 & 話そうとする試みが毎分約$60$語の速さでテキストに変換される & 明瞭な発話を失った人、皮質に配列：毎分$62$語。語彙$50$語で単語誤り率$9.1\,\%$、語彙$125\,000$語で$23.8\,\%$。 & \cite{Willett2023} & 測定済み \\
12 & インターフェースが取り出すのは容量のごく一部：1ニューロンで毎秒数十ビット、100個で数千ビット（命題~\ref{prop:probe}） & 上限~\eqref{eq:spikeentropy}は理論。8行目、10行目との比較は本章の結論。ボトルネックが配線ではなく低次元性と符号の無知であることは、直接の実験で確かめられていない。 & \cite{Rieke1997} & 理論 \\
13 & 脳とデコーダは互いから学ぶ & サルがカーソルを操作。デコーダは閉ループで適応し、ニューロンも同時に学習し直す。精度は上がり、技能は保持され、自分の腕の動きの影響を受けにくくなる。学習速度を分けるという助言は工学上の経験則であり、本章に個別の実験はない。 & \cite{Orsborn2014coad} & 測定済み \\
14 & 電極を流れる電流は、周囲のニューロンと配線を種類の区別なく興奮させる & マウスとネコの皮質、二光子カルシウムイメージング：弱い電流でも、直径数十マイクロメートルの体積内の軸索を直接興奮させることで、最大数ミリメートルに散らばるまばらなニューロンを興奮させる。つまり興奮は非選択的だが、一様ではない。 & \cite{Histed2009stim} & 測定済み \\
15 & 視覚皮質の刺激は点を灯す。同時に最大$16$点で一部の文字と物体の境界が認識できる & 盲目の人、視覚皮質に$96$電極を$6$か月：単一電極の閾値は$66.8\pm36.5$~\textmu A。電極群の同時刺激（最大$16$電極は刺激装置の上限）で閾値が下がり、区別できるパターンが生じ、一部の文字と物体の境界が認識される。 & \cite{Fernandez2021} & 測定済み \\
16 & Neuralinkの第二の方向は、視覚皮質に信号を入れる装置 & 開発についての同社の発表のみ。その動作に関する査読付きデータは見つからなかった。 & Neuralinkの報告 & 仮説 \\
17 & ブロードキャスト伝達物質の濃度は組織内の化学センサで測れる & ラット脳スライス、炭素繊維による高速サイクリックボルタンメトリー：セロトニンの放出と再取り込みを測定。物質はシナプスの外へ出る。 & \cite{Bunin1998} & 測定済み \\
\bottomrule
\end{longtable}}

\section{第\ref{sec:eetypes}章　素子としてのニューロン}

個々の細胞の動作モードは、電極からの電流注入と電圧記録で直接測定されている。積分器とクラス分けの数学は古典的な理論である。脳がモードの切り替えを計算に使っているというのは、依然として仮説である。

{\footnotesize
\setlength{\tabcolsep}{3pt}\renewcommand{\arraystretch}{1.15}
\begin{longtable}{>{\raggedright}p{0.5cm}>{\raggedright}p{4.0cm}>{\raggedright}p{5.6cm}>{\raggedright}p{2.3cm}>{\raggedright\arraybackslash}p{1.7cm}}
\toprule
№ & 主張 & 実験と測定内容 & 出典 & 状態 \\
\midrule\endfirsthead
\toprule
№ & 主張 & 実験と測定内容 & 出典 & 状態 \\
\midrule\endhead
1 & 共振器：電圧の変化に逆らう遅い電流が、ニューロンを数Hzから数十Hzに最大をもつバンドパスフィルタにする（第\ref{sec:resonator}節） & スライス中の\nt{GLUT}ニューロンに周波数が滑らかに上がる正弦波電流を注入し、インピーダンス$|Z(f)|$を測定：最大は$33^\circ$Cで$2$--$5$~Hz（$38^\circ$Cで約$7$~Hz）。共振は遅いカリウム電流と陽イオン電流が生み、持続性ナトリウム電流が強める。総説は多くの細胞クラスで同じ手法を示す & \cite{HuVervaekeStorm2002,HutcheonYarom2000} & 測定済み \\
2 & 遅い電流はインダクタンスのように振る舞い、膜容量と合わせて共振回路を作る & 電流に対する軸索の閾値下応答を測定し、線形化したコンダクタンス方程式で記述。等価回路は電圧に依存する「現象論的」インダクタンスを含む & \cite{Mauro1970} & 理論 \\
3 & フィードバックのある群の時定数は$\tau_{\text{eff}}=\tau/(1-w)$~\eqref{eq:perfectint}。$\tau=0.1$~s、$w=0.995$で$20$~s & 群の線形方程式から本章で導出。眼の位置を保持する機構として理論研究で提案された & 本章で導出；\cite{Seung1996} & 理論 \\
4 & 眼の位置の積分器は、個々の細胞の性質ではなくネットワークで入力を保持する & 眼の位置を保持する魚の神経核での細胞内記録：急速な回転のあと、ニューロンの発火率は階段状に変わって保たれる。1個の細胞への短い電流は一時的なずれしか起こさず、電圧の階段は閾値下でも保たれる。それを維持するのはシナプス入力である & \cite{Aksay2001} & 測定済み \\
5 & 漏れのない積分器には学習による絶え間ない調整が必要。調整が狂うと忘れるか飽和する & 眼の位置の$\pm$に比例する視覚フィードバックで魚を訓練：保持は不安定または漏れのあるものになり、ニューロンの時定数は1桁以上下がって$<1$~sになった。通常の視覚フィードバックで安定性は徐々に戻った & \cite{Major2004} & 測定済み \\
6 & 双安定ニューロン：短い入力が長いプラトーをオンにし、抑制がオフにする。この性質はセロトニンでオンになる（第\ref{sec:bistable}節） & ネコの筋肉を制御する\nt{ACET}ニューロンの細胞内記録：短いシナプス興奮でニューロンは長く続く発火に移り、短い抑制でリセットされる。脊髄ネコではセロトニン前駆体の投与後にこの状態が現れる & \cite{Hounsgaard1988} & 測定済み \\
7 & 二つのクラス：積分器型（発火率がゼロから上がる）と共振器型（閾値ですぐ有限の発火率） & クラスは分岐理論から導かれる。実験：皮質スライスで\nt{GLUT}ニューロンはいくらでも低い発火率で発火し始め（タイプ1）、速い\nt{GABA}ニューロンは有限の発火率から跳躍的に始まる（タイプ2） & \cite{Izhikevich2007,Tateno2004} & 測定済み \\
8 & 1個のニューロンが積分器にも同時検出器にもなる：抑制が加算の窓を短くする & 海馬スライスで、興奮性入力が閾値まで加算される窓は$2$~ms未満。興奮の直後に届くフィードフォワード抑制がこれを短くする。抑制の弱い樹状突起では窓が広い & \cite{Pouille2001} & 測定済み \\
9 & 軸索による遅延線（表~\ref{tab:eetypes}） & メンフクロウで、同時検出器へ向かう軸索の遅延を測定：軸索の長さの違いが異なる遅延を生み、検出器は音の到達時間差に同調している & \cite{Carr1990} & 測定済み \\
10 & 覚醒中はペースメーカー、睡眠中は黙った細胞 & \nt{SERO}ニューロンは日中ほぼ規則的に発火し、徐波睡眠で遅くなり、レム睡眠ではほぼ沈黙する（詳しくは第\ref{sec:sleep}章） & \cite{JacobsAzmitia1992,McGintyHarper1976} & 測定済み \\
11 & 素子を決めるのはイオンチャネルと、それを調整するブロードキャストチャネル（命題~\ref{prop:eetypes}） & 小さなネットワークで示された：調節物質がニューロン固有の性質とシナプスの強さを変え、同じ回路が異なるリズムを出す & \cite{Marder2012} & 測定済み \\
12 & 脳はモードの切り替えを計算に使っている（命題~\ref{prop:eetypes}） & 細胞のモードの切り替えが課題の解決に必要であることを示す直接の実験はない。調節物質が回路の出力を変える実験があるだけである & \cite{Marder2012} & 仮説 \\
\bottomrule
\end{longtable}}

\section{第\ref{sec:dendrites}章　樹状突起の数}

クラスの構造と入力の数は、解剖学的方法と電気的記録で測定されている。見積もり~\eqref{eq:dendriteload}はコンデンサの単純な物理である。入力数の計算上の説明は、理論とモデルに依拠している。

{\footnotesize
\setlength{\tabcolsep}{3pt}\renewcommand{\arraystretch}{1.15}
\begin{longtable}{>{\raggedright}p{0.5cm}>{\raggedright}p{4.0cm}>{\raggedright}p{5.6cm}>{\raggedright}p{2.3cm}>{\raggedright\arraybackslash}p{1.7cm}}
\toprule
№ & 主張 & 実験と測定内容 & 出典 & 状態 \\
\midrule\endfirsthead
\toprule
№ & 主張 & 実験と測定内容 & 出典 & 状態 \\
\midrule\endhead
1 & 膜の比容量は$c_m\approx1$~\textmu F/cm$^2$ & ニューロンの細胞体から膜をピペットに取り、電圧ステップを与えて容量電流の減衰から全容量を求め、面積で割った：3クラスのニューロンで$0.9$~\textmu F/cm$^2$ & \cite{Gentet2000} & 測定済み \\
2 & 充電時間$t\approx c_mA\Delta V/I$~\eqref{eq:dendriteload}：大きなニューロンには数十個の同時入力が必要で、小さなニューロンには大きなシナプス1個で足りる & コンデンサの充電法則による見積もり。数値（$20$と$300$~pF、$0.1$と数nA）は桁の値 & 本章で導出 & 理論 \\
3 & 1個の巨大シナプスは数nAの電流を流し、小さなニューロンをすぐに閾値に届かせる & スライスで終末と受け手を同時記録：シナプス1個の電流は$2$--$13$~nA、入力スパイクのたびに閾値を超える応答が生じ、遅れは$36^\circ$Cで約$0.4$~ms。受け手は\nt{GLYC}を出す & \cite{Borst1995,BorstSoria2012} & 測定済み \\
4 & 樹状突起ゼロ：触覚と痛みの感覚ニューロンでは、スパイクはセンサから細胞体を経ずに脊髄へ向かう & 構造（細胞体の近くで二つに分かれる軸索）は解剖学的に確立。伝導に細胞体の興奮性が要らないことは、このニューロンの検証済みモデルで示された：興奮性のない細胞体でも、スパイクは分岐点を同じく確実に通過する & \cite{AmirDevor2003} & 理論 \\
5 & 樹状突起1本：嗅覚ニューロンは1種類の受容体をもち、1本の入力線を伝える & 受容体遺伝子の実験の総説：各嗅覚ニューロンは大きな遺伝子ファミリーのうち一つの受容体遺伝子を発現する & \cite{Mombaerts2004} & 測定済み \\
6 & 樹状突起2本：音の方向検出器では、左耳と右耳からの入力が別々の樹状突起に接続する & 総説にまとめられた哺乳類の解剖と記録：両耳からの入力は反対向きの2本の樹状突起に分かれている & \cite{Grothe2010} & 測定済み \\
7 & 入力を2本の樹状突起に分けるとANDが厳密になる & モデルで示された：1本の樹状突起での飽和~\eqref{eq:shunt}により、片側からの入力は閾値に届かず、同時検出が改善する。入力の配置を変えた直接の実験はない & \cite{AgmonSnir1998} & 理論 \\
8 & 運動学習回路の小さな\nt{GLUT}ニューロンは約$7\cdot10^{10}$個で、それぞれ$\sim4$個の入力 & 等方性分画法による細胞計数：ヒトの脳のこの部分に約$690$億個のニューロン。生きたラットでの単一細胞記録：接触に対して少数の入力から離散的な電流のバーストが届き、それらが加算されると細胞が発火する & \cite{Azevedo2009,Chadderton2004} & 測定済み \\
9 & $n$入力の閾値素子が分けられるランダムなパターンは$P_{\max}\approx2n$個まで~\eqref{eq:cover}。運動学習回路の出力ニューロンでは上限$\sim2\cdot10^5$、現実的な見積もりは$\sim4\cdot10^4$個の対応 & 上限は組合せ幾何学の古典的な結果。現実的な見積もりは、正の重みのみとノイズへの余裕を仮定した容量理論を、このニューロンの測定された入力数に適用したもの & \cite{Cover1965,Brunel2004} & 理論 \\
10 & 入力の少ないミキサは入力の拡張に必要で、「4」は最適に近い & 理論：ミキサ層が与える表現の次元は、解剖学的に観察される入力数付近で最大になる。拡張の元の仮説は古典的な研究にある & \cite{Marr1969,Albus1971,LitwinKumar2017} & 仮説 \\
11 & 大きな木：$10^4$--$10^5$個の入力 & 電子顕微鏡像によるシナプスの計数：ラットの運動学習回路の出力\nt{GABA}ニューロン1個に$\approx1.7\cdot10^5$個のシナプス。海馬の\nt{GLUT}ニューロンに約$30\,000$個の興奮性入力と$1\,700$個の抑制性入力 & \cite{NapperHarvey1988,Megias2001} & 測定済み \\
12 & 大きな木の枝は、それぞれ閾値をもつ別々のサブ回路であり、ニューロンは小さな多層ネットワークである & 細い樹状突起の2か所を局所刺激：同じ枝の入力はシグモイド状に、異なる枝の入力は線形に加算される。ヒト皮質ニューロンの樹状突起で段階的なカルシウムスパイクが見つかり、1個の細胞が線形分離不可能な課題を解ける。モデル：1個のニューロンを再現するには深いネットワークが必要 & \cite{Polsky2004,Gidon2020,Beniaguev2021} & 測定済み \\
13 & 出力を兼ねる樹状突起：網膜の\nt{GABA}ニューロンの各枝は独立した方向検出器 & 個々の枝での二光子カルシウム記録：枝の信号は細胞体から先端への動きに選択的だが、細胞体の電圧はそうではない & \cite{Euler2002} & 測定済み \\
14 & 入力数は課題に従う（命題~\ref{prop:dendrites}） & クラスの構造は測定済み（3--13行目）。入力数が速さや分類のために選ばれたことは、個々のクラスについて理論とモデルで示されているだけである & \cite{LitwinKumar2017,AgmonSnir1998} & 仮説 \\
\bottomrule
\end{longtable}}

\section{第~\ref{sec:sleep}~章　睡眠}

睡眠のモード、リプレイ、シナプスの縮小は、ニューロンの記録、マイクロダイアリシス、電子顕微鏡で測定されている。睡眠のスケジュールと遅い振動は本書のモデルで記述され、睡眠の目的は主に仮説である。

{\footnotesize
\setlength{\tabcolsep}{3pt}\renewcommand{\arraystretch}{1.15}
\begin{longtable}{>{\raggedright}p{0.5cm}>{\raggedright}p{4.0cm}>{\raggedright}p{5.6cm}>{\raggedright}p{2.3cm}>{\raggedright\arraybackslash}p{1.7cm}}
\toprule
№ & 主張 & 実験と測定内容 & 出典 & 状態 \\
\midrule\endfirsthead
\toprule
№ & 主張 & 実験と測定内容 & 出典 & 状態 \\
\midrule\endhead
1 & 睡眠中も皮質のニューロンは沈黙しない。変わるのは動作モードである & ラット皮質ニューロンの長時間記録：ノンレム睡眠でも発火率はゼロにならず、活動期と静止期が交互に現れる。長い覚醒のあとはすべての状態で発火率が高く、睡眠のあとは低い & \cite{Vyazovskiy2009} & 測定済み \\
2 & \nt{ACET}は日中とレム睡眠で高く、ノンレム睡眠で低い（表~\ref{tab:sleepmodes}） & 自由に動くネコの皮質でのマイクロダイアリシス：覚醒時のアセチルコリン放出はノンレム睡眠より$100$--$175\%$高く、レム睡眠ではおよそ安静覚醒のレベル & \cite{Marrosu1995} & 測定済み \\
3 & \nt{NORA}と\nt{SERO}はノンレム睡眠で静まり、レム睡眠ではほぼ沈黙する & ラットの\nt{NORA}ニューロンとネコの\nt{SERO}ニューロンを睡眠段階ごとに単一記録：発火率は覚醒時に最大、ノンレム睡眠で低く、レム睡眠ではほぼゼロ & \cite{AstonJonesBloom1981,McGintyHarper1976} & 測定済み \\
4 & レム睡眠では筋肉が\nt{GABA}と\nt{GLYC}による抑制で切り離される & ラットで咀嚼筋の\nt{ACET}ニューロンの活動を測り、そこへ直接遮断薬を与えた：麻痺が解けるのは、両タイプの\nt{GABA}受容体と\nt{GLYC}受容体をすべて遮断した場合だけ & \cite{BrooksPeever2012} & 測定済み \\
5 & 睡眠圧$S$は二つの閾値をもつコンデンサ~\eqref{eq:sleeppressure}、$\tau_r=18.2$~h、$\tau_d=4.2$~h & 2過程モデル。時定数は脳波の徐波パワーが覚醒とともに増え睡眠中に減る様子から、閾値の形は自発的な覚醒の時刻から得られた & \cite{Borbely1982,Daan1984} & 理論 \\
6 & マシンはおのずと$16.1$~hの覚醒と$8.0$~hの睡眠に落ち着く（図~\ref{fig:twoprocess}） & 揺れる閾値を用いた\eqref{eq:sleeppressure}による計算、スクリプト\texttt{models/sleep\_models.py}：覚醒$16.1$~h、睡眠$7.9$--$8.0$~h & --- & モデル \\
7 & $40$~h眠らないと$S=0.90$になるが、睡眠は$8.8$~hしか続かない。借りは長さではなく深さで返される & モデル：\texttt{models/sleep\_models.py}は$S=0.90$と$8.8$~hを与える。実験：$40.5$~hの断眠後、8人の被験者で回復第1夜の脳波の徐波パワーが通常より有意に高い & \cite{Borbely1981} & モデル \\
8 & 物理的には$S$はアデノシンである & ネコでのマイクロダイアリシス：覚醒を制御する領域の一つで細胞外アデノシンが覚醒中に増え、断眠でさらに増え、睡眠中に減る。$S$がアデノシンだけであることは示されていない & \cite{PorkkaHeiskanen1997,Dunwiddie2001} & 仮説 \\
9 & ノンレム睡眠で皮質は振動する：数百msの活動と静止、1秒に1回未満（$<1$~Hz、麻酔下では多くが$0.3$--$0.4$~Hz） & 麻酔下のネコでの細胞内記録：$0.8$--$1.5$~sの脱分極と長い過分極、リズム$<1$~Hz、最多は$0.3$--$0.4$~Hz。同じ活動と静止の交代は自然睡眠でも見られる（1行目） & \cite{Steriade1993,Vyazovskiy2009} & 測定済み \\
10 & 振動はフィードバックと疲労をもつ群~\eqref{eq:updown}。$b=5$で周期$1.1$~s（モデルの値）、活動は時間の約$35\%$。$b=1$では一様な活動（図~\ref{fig:updown}） & 計算、スクリプト\texttt{models/sleep\_models.py}：周期$1.11$~s、活動時間の割合$0.35$（整数個の周期での平均）。$b=1$では活動の割合$1.0$。モデルの周期は測定値（9行目）より短い & --- & モデル \\
11 & \nt{ACET}は振動を止め、皮質を一様な活動に移す & ネコで\nt{ACET}ニューロンの核を刺激すると、皮質ニューロンの$79\%$で遅いリズムが遮断され、主に長い過分極が消えることで一様な活動に移った。アセチルコリンは疲労を生む遅いカリウム電流を抑える & \cite{SteriadeAmzica1993,McCormick1992} & 測定済み \\
12 & $7$--$15$~Hzのリズムのバーストは、皮質と中継核の間の\nt{GLUT}--\nt{GABA}ループが生む & フェレットの視覚中継核のスライスで、バーストは、中継細胞自身が興奮させる隣接\nt{GABA}核からの抑制に対して中継細胞が返すバーストによって生じる & \cite{vonKrosigk1993,SteriadeMcCormickSejnowski1993} & 測定済み \\
13 & 1日のリプレイは早送りで行われる：海馬ではおよそ$20$倍、皮質では$6$--$7$倍速い & ノンレム睡眠で、海馬ニューロンの繰り返す系列が時間的に圧縮されて現れる。トラックを走ったあと、場所ニューロンの系列が同じ順序で$\sim100$~msの短いバーストとして、およそ$20$倍に圧縮されて再生された。皮質での圧縮は$6$--$7$倍 & \cite{Nadasdy1999,LeeWilson2002,Euston2007} & 測定済み \\
14 & リプレイと皮質の協調を強めると記憶が良くなる（因果関係） & ラットで学習後、リプレイに同期した刺激によって皮質の徐波とリズムのバーストとの結びつきを強めた：翌日の記憶は高く、対照群では偶然レベル & \cite{Maingret2016} & 測定済み \\
15 & リプレイの目的は遅い記憶のインターリーブ学習 & 二つの学習系の理論と人工ネットワークでの計算：逐次学習は古いものを壊し、インターリーブ学習は壊さない。リプレイがまさにこのために必要だという脳での直接の実験はない & \cite{McClelland1995} & 仮説 \\
16 & 睡眠後、シナプスは大きさに比例して縮み、大きなものはほとんど変わらない & マウス皮質の$6920$個のシナプスの3次元電子顕微鏡観察：接触面積は睡眠後に覚醒後より$\sim18\%$小さく、縮小は大きさに比例し、大きく安定したシナプスは変わらない & \cite{deVivo2017} & 測定済み \\
17 & 睡眠の主な役割は重みの再正規化 & シナプス恒常性仮説。1行目と16行目はこれと整合するが、それが主な機能であることは証明しない & \cite{TononiCirelli2014} & 仮説 \\
18 & レム睡眠は世界モデルのベンチ試験 & モード（表~\ref{tab:sleepmodes}）は測定済み。ネットワークが世界モデルを試運転しているというのは理論上の推測で、実験はない & \cite{Hobson2009} & 仮説 \\
19 & 睡眠中の脳の洗浄：未解決の問題 & ある実験：睡眠中は細胞間隙が$60\%$広く、除去が速い。測定方法の異なる別の実験：睡眠中と麻酔下では除去が明らかに遅い。結果は互いに矛盾する & \cite{Xie2013,Miao2024} & 仮説 \\
20 & 局所睡眠：覚醒している動物の皮質の部位が短時間静止に落ちる & 長い覚醒のあとのラットで、皮質の一部位のニューロンがノンレム睡眠のように短く沈黙し、局所脳波に徐波が現れる。その瞬間、ラットは課題で誤りやすい & \cite{Vyazovskiy2011} & 測定済み \\
\bottomrule
\end{longtable}}

\section{第\ref{sec:livingmachine}章　生きたニューロンでできたマシン}

電極アレイ上の培養の振る舞いは、記録と刺激を伴う直接の実験で測定されている。一斉発火の機構はシナプス枯渇のモデルと微小培養での実験に基づき、シャーレの中のドーパミンに関する主張は単発の実験に基づく。そのうちいくつかは著者ら自身がデモンストレーションにすぎないとみなしている。

{\footnotesize
\setlength{\tabcolsep}{3pt}\renewcommand{\arraystretch}{1.15}
\begin{longtable}{>{\raggedright}p{0.5cm}>{\raggedright}p{4.0cm}>{\raggedright}p{5.6cm}>{\raggedright}p{2.3cm}>{\raggedright\arraybackslash}p{1.7cm}}
\toprule
No. & 主張 & 実験と測定されたもの & 出典 & 状態 \\
\midrule\endfirsthead
\toprule
No. & 主張 & 実験と測定されたもの & 出典 & 状態 \\
\midrule\endhead
1 & チップ上の培養：大部分は\nt{GLUT}、少数の\nt{GABA}の割合は標本による。海馬の培養では約 $6\%$ & 電極アレイ上の数千ニューロンのネットワークは標準的な実験である（行~3）。海馬の培養での\nt{GABA}ニューロンの割合は GABA 染色で数えられ、ニューロンの約 $6\%$ だった。大脳皮質の培養については検証された数値を挙げない & \cite{Benson1994} & 測定済み \\
2 & オルガノイド：ヒトの幹細胞からできた、大きさ約1ミリの三次元神経組織 & 幹細胞から、複数の脳領域をもつ三次元オルガノイドが育てられた。その中には前駆細胞の層と成熟ニューロンをもつ大脳皮質も含まれる & \cite{Lancaster2013} & 測定済み \\
3 & 入力がないと培養は、培養によって1秒から数分の間隔で一斉発火する & 密度 $3\,000$ から $50\,000$ ニューロンの大脳皮質培養 $58$ 個を5週間記録した（30分の記録 $963$ 本）。2週間後にはほぼすべての培養でネットワーク全体の一斉発火が優勢になる。一斉発火の間隔は $1$ から $300$~s で、標本に強く依存する & \cite{Wagenaar2006} & 測定済み \\
4 & 一斉発火は伝達物質の枯渇で終わり、間隔は $T_{\text{burst}}\approx\tau_{\text{rec}}\ln\frac{1}{1-x^*}$~\eqref{eq:burstinterval}。$2$--$4$~s は $\tau_{\text{rec}}=0.8$~s のモデルによる推定で、測定された回復はもっと遅い & 式は本章で導いた。モデル：枯渇するシナプスをもつネットワークは、自らネットワーク全体の一斉発火を生む。$4$--$30$ ニューロンの微小培養での実験：一斉発火の持続時間は前回からの時間に依存し、伝達物質の小胞を枯渇させると一斉発火は消える。著者らの結論は、一斉発火は伝達物質の蓄えで制限されるというもの。そこでの蓄えの回復は数十秒で、同じ式で数十秒から数分の間隔を与える & 本章の導出、\cite{Tsodyks2000,CohenSegal2011,Tsodyks1997} & 理論 \\
5 & 「黙る教師」：目的の応答が出たら刺激を止め、応答が定着する & 培養を $0.3$--$1$~Hz で刺激し、刺激後 $50\pm10$~ms に選んだ応答が現れたら刺激を $5$~min 止めた。数サイクル後、目的の応答はすぐに誘発されるようになった。学習曲線が描かれた & \cite{Shahaf2001} & 測定済み \\
6 & 培養がテニスをする。セッション内のラリーの有意な伸びは完全なフィードバックのときだけ & 詳細は第\ref{sec:dishbrain}章とその補遺。刺激の代わりに無音にしても、セッションの終わりにはフィードバックなしより上手になるが、効果は小さい。セッション間の学習は不安定 & \cite{Kagan2022} & 測定済み \\
7 & 培養が自分の入力を予測するように学習するというのは仮説。「知性」をめぐる大げさな言葉には反論がある & 自由エネルギー原理は理論的な提案であり、より単純な説明と区別する直接の実験はない。30人の研究者が反論論文で、ループ内での振る舞いの変化は知性も感覚も示さないと指摘した & \cite{Friston2010,Balci2023} & 仮説 \\
8 & 刺激パターンを選べば、培養は仮想の体を目標へ動かす & プログラムが現在の成績に応じて訓練刺激のパターンを自動で選んだ。数十分でネットワークは指定した状態に達し、$2$~h の訓練後、可塑性は $80$~min 続いた。振る舞いと無関係に与えた同じ刺激は効果がなかった & \cite{Bakkum2008} & 測定済み \\
9 & リザバー：学習するのは線形読み出し $y=W_{\text{out}}\mathbf r$~\eqref{eq:reservoir} だけ & リザバー計算の理論：減衰する記憶をもつ任意の非線形システムに、学習させた線形出力を加える & \cite{Maass2002,JaegerHaas2004} & 理論 \\
10 & リザバーとしてのオルガノイドは話者を約 $78\%$ の精度で区別する（著者らのデータ）。誤差で学習するのは読み出しで、オルガノイドの結合は教師なしで変わる & 複数の話者の音声を、電極アレイ上の電流パターンとしてオルガノイドに与えた。学習させた読み出しは約 $78\%$ の精度で話者を区別した。これは著者らの報告である。著者らはまた、訓練中にオルガノイドの機能的結合が変化したと報告し、これをオルガノイド内の教師なし学習と呼んでいる & \cite{Cai2023} & 測定済み \\
11 & 100万個のニューロンがスパイクに使う電力は約 $10^{-4}$~W~\eqref{eq:dishpower} & 推定：大脳皮質のエネルギー収支\eqref{eq:energy}から得たスパイクのコスト $10^{-10}$~J にスパイク数を掛けた。培養の電力の直接測定はない & 本章の導出、\cite{Attwell2001} & 理論 \\
12 & 光のフラッシュで放出するドーパミン：ケージドドーパミン $1$~mM、$240$ 回の繰り返しで誘発スパイク数が増加 & オルガノイドの遠隔プラットフォームで、刺激がより多くのスパイクを誘発したとき、光感受性ケージに入ったドーパミン（$1$~mM）を紫外線（$365$~nm、$0.8$~s）で解放した。$240$ ステップ（約 $5$~h）でスパイク数は全体として増えた。著者らは、実験一つではドーパミンとの関係は確立できないと書いている。対照はない & \cite{Jordan2024} & 仮説 \\
13 & 生きた細胞からのドーパミンが有機トランジスタの重みを長期かつ可逆的に変える & ドーパミンを分泌する細胞を有機ニューロモルフィック素子の上で育てた。電極でのドーパミンの酸化がチャネルのコンダクタンスを変える。重みの長期変化とその復帰が示された & \cite{Keene2020} & 測定済み \\
14 & 働く\nt{DOPA}ニューロンをもつオルガノイドは存在する。そのドーパミンを教師として使った例はない & ヒトの幹細胞からのオルガノイドに、電気的に活動する成熟\nt{DOPA}ニューロンとドーパミン産生が見つかった。このドーパミンが別のネットワークを教える実験は文献に見つからなかった & \cite{Jo2016} & 測定済み \\
15 & オルガノイドが棒を立てる。プログラムによる学習刺激はランダムより良いが定着せず、\nt{GLUT}シナプスを介する & マウス大脳皮質のオルガノイドを「台車の上の棒」課題と閉ループでつないだ。大多数のオルガノイドで、強化学習が選んだ信号はランダムな信号や信号なしより良い成績をもたらした。$45$~min の休止後、改善は残らなかった。AMPA 受容体と NMDA 受容体の遮断で改善は消えた & \cite{Robbins2026} & 測定済み \\
16 & 制御レジスタがなければ、テストベンチ上のネットワークは何を成功とみなすかを自分で決められない（命題\ref{prop:livingmachine}） & 行 5--15 のすべての実験で、学習信号は実験者かプログラムが与える。培養が自ら成功の基準を作り出した実験はない。これは不在からの推論であり、実験ではない & --- & 仮説 \\
\bottomrule
\end{longtable}}

\section{第\ref{sec:dishbrain}章　DishBrain}

テストベンチの構成とゲームの結果は、一つの実験研究とその続編から採った。ノイズと良い設定への漂流の式は本章で導き、本書のモデルで確かめた。シャーレの中の学習の仕組みは確定していない。

{\footnotesize
\setlength{\tabcolsep}{3pt}\renewcommand{\arraystretch}{1.15}
\begin{longtable}{>{\raggedright}p{0.5cm}>{\raggedright}p{4.0cm}>{\raggedright}p{5.6cm}>{\raggedright}p{2.3cm}>{\raggedright\arraybackslash}p{1.7cm}}
\toprule
No. & 主張 & 実験と測定されたもの & 出典 & 状態 \\
\midrule\endfirsthead
\toprule
No. & 主張 & 実験と測定されたもの & 出典 & 状態 \\
\midrule\endhead
1 & テストベンチ：チップあたりマウス大脳皮質の細胞約 $800\,000$ 個またはヒト細胞約 $10^6$ 個。$8$~mm$^2$ に電極 $26\,000$ 本、記録は最大 $1024$、刺激は $8$ 本 & 論文の方法：MaxOne チップ、$8$~mm$^2$ に白金電極 $26\,000$ 本、$1024$ チャネルを毎秒 $20\,000$ サンプルで記録。技術的には最大 $32$ 電極を刺激できるが、独立に刺激できたのは $8$ 本。マウス培養はおよそ $10$ 日目から使用 & \cite{Kagan2022} & 測定済み \\
2 & $10$~ms 周期のループ：運動領域のスパイクがラケットを動かす（図\ref{fig:dishloop}） & スパイクは $10$~ms の窓ごとに数えられる。スパイクから応答刺激までの遅れは約 $5$~ms で、装置のバッファで決まる & \cite{Kagan2022} & 測定済み \\
3 & 入力：場所による符号（$8$ 本のどの電極か）と頻度 $4$--$40$~Hz & 主実験では刺激頻度が、ボールが遠い壁にあるときの $4$~Hz からラケットのそばの $40$~Hz まで線形に増える。ボールのパルスの振幅は $75$~mV、パルスは矩形の二相性 & \cite{Kagan2022} & 測定済み \\
4 & 出力 $\Delta p=c\,(n_\uparrow-n_\downarrow)$~\eqref{eq:dishreadout}、窓あたりのカウントのノイズは $1/\sqrt{RT}$ & 二つの領域の差をとる規則は論文に記述がある（活動に応じた領域の重みつき）が、正確な式はない。ノイズの推定はポアソン計数からの帰結で、本章で導いた & \cite{Kagan2022}、本章の導出 & 理論 \\
5 & 教師：ヒットの後に予測可能な刺激 $75$~mV、$100$~Hz、$0.1$~s。ミスの後に予測不能な刺激 $150$~mV、$5$~Hz、$4$~s をランダムな場所と時刻に、その後 $4$~s 休止 & 方法：ヒットの後、$75$~mV、$100$~Hz、$100$~ms を $8$ 本の電極すべてに同時に。ミスの後、$150$~mV、$5$~Hz をランダムな電極にランダムな時刻で $4$~s、その後 $4$~s 刺激なし。培養にドーパミンはない & \cite{Kagan2022} & 測定済み \\
6 & ネットワークは自分の入力が予測可能になるように行動する & 自由エネルギー原理は理論的な提案で、著者らはそこからプロトコルを導いた。実験はこれと矛盾しないが、より単純な仕組み（行~7--8）と区別できない & \cite{Friston2010,Kagan2022} & 仮説 \\
7 & 失敗がネットワークを揺さぶり成功が揺さぶらないなら、設定で過ごす時間は $\rho\propto1/P_{\text{miss}}$~\eqref{eq:dishdensity}（命題\ref{prop:dishscramble}） & 対称な揺さぶりをもつマルコフ連鎖の流れの釣り合いから従う。本章で導いた。培養での直接の検証はない & 本章の導出 & 理論 \\
8 & 「ミスでかき混ぜる」モデルは学習する：ヒットの割合 $0.21\to0.38$。毎ラリー後に揺さぶると $\approx0.12$、揺さぶりなしで $0.19$（図\ref{fig:dishmodel}） & 計算、スクリプト \texttt{models/dishbrain\_model.py}、モデル培養 $40$ 個に各 $2000$ ラリー：$0.21\to0.38$、$0.13\to0.11$、$0.19\to0.19$ & --- & モデル \\
9 & ラリーが長くなるのは完全なループの生きた培養だけ。対照は学習しない & $20$~min のセッション $399$ 回：マウス培養 $9$ 個（$101$ セッション）、ヒト $11$ 個（$138$）、培地のみのチップ $6$ 枚（$80$）、安静の培養 $20$ 個（$42$）、シミュレーション $3$ 個（$38$）。最後の $15$~min の平均ラリー長が最初の $5$~min より長いのは生きた培養だけ。ニューロンでない細胞（HEK293T）と培地のみのチップでは効果なし & \cite{Kagan2022} & 測定済み \\
10 & セッション内の有意な伸びは完全なフィードバックのときだけ。刺激の代わりに無音でも、セッションの終わりにはフィードバックなしより上手だが、効果は小さい & $3$ 日で $486$ セッション：「刺激」（$n=27$）、「無音」（$17$）、「フィードバックなし」（$15$）。時間による伸びが有意なのは「刺激」条件だけ。セッションの終わりに「無音」条件は「フィードバックなし」を小さな効果で上回った & \cite{Kagan2022} & 測定済み \\
11 & 効果は小さい。ネットワークが長期的に学習するかは未解決 & セッション内の改善は確かだが、著者ら自身によれば、数日にわたるセッション間の学習は安定しては観察されなかった & \cite{Kagan2022} & 測定済み \\
12 & プレイ中、ネットワークは臨界状態に近づき、近いほどプレイが良い & 同じ培養の活動の雪崩の解析：構造化された入力はネットワークを臨界状態に近づけ、プレイの質はそこへの近さと相関する。ただし、行動の結果についての情報がなければ、臨界性だけでは学習に足りない & \cite{Habibollahi2023} & 測定済み \\
13 & 培養の「知性」は示されていない & 30人の研究者による反論論文：閉ループでの振る舞いの変化は知性も感覚も証明しない。それを示す実験はない & \cite{Balci2023} & 仮説 \\
14 & 提案する四つ目の対照：ミスの後に同じ長さとエネルギーの予測可能な刺激（\ref{sec:dishspec}節） & この実験は行われていない。本書の提案である & --- & 仮説 \\
\bottomrule
\end{longtable}}

